\documentclass[journal]{IEEEtran}
\usepackage[T1]{fontenc}
\usepackage{algorithmic}
\usepackage{algorithm}
\usepackage{array}
\usepackage{textcomp}
\usepackage{stfloats}
\usepackage{url}
\usepackage{verbatim}
\usepackage{graphicx}
\usepackage{amsfonts}
\usepackage{amsmath}
\usepackage{amssymb}
\usepackage{amsthm}
\theoremstyle{plain}
\newtheorem{dfn}{Definition}
\newtheorem{theo}{Theorem}
\newtheorem{rem}{Remark}
\usepackage{color}
\usepackage{booktabs}
\usepackage{multirow}
\usepackage{cite}
\makeatletter
\let\MYcaption\@makecaption
\makeatother
\usepackage[justification=centering,labelfont=sf,textfont=sf]{subcaption}
\makeatletter
\let\@makecaption\MYcaption
\makeatother
\usepackage{setspace}
\usepackage{cellspace}
\usepackage{newcommands}
\usepackage{hyperref}
\hypersetup{
    colorlinks = true, % リンクに色をつける
    urlcolor = black, % URLの色
    linkcolor = black, % 内部リンクの色
    citecolor = black, % 引用リンクの色
    anchorcolor = black, % アンカーテキストの色
    pdfborder = {0 0 0} % リンクの枠を消す
}
\usepackage{enumitem}

\hyphenation{op-tical net-works semi-conduc-tor IEEE-Xplore}

\begin{document}
\title{Generalized Graph Signal Sampling by Difference-of-Convex Optimization}

\author{Keitaro~Yamashita,~\IEEEmembership{Student Member,~IEEE,}
        Kazuki~Naganuma,~\IEEEmembership{Member,~IEEE,}
        and~Shunsuke~Ono,~\IEEEmembership{Senior Member,~IEEE}
\thanks{This work was supported in part by JST PRESTO under Grant JPMJPR21C4, JST AdCORP under Grant JPMJKB-2307, JST ACT-X Grant JPMJAX23CJ, JSPS KAKENHI under Grant 22H03610, 22H00512, 23H01415, 23K17461, 24K03119, and 24K22291, and Grant-in-Aid for JSPS Fellows under Grant 23KJ0912.}
\thanks{K. Yamashita, K. Naganuma, and S. Ono are with the Department
of Computer Science, Institute of Science Tokyo, Tokyo,
Japan e-mail: yamashita.k.45d1@m.isct.ac.jp; naganuma.k.491a@m.isct.ac.jp; ono.s.5af2@m.isct.ac.jp (see https://www.mdi.c.titech.ac.jp/).}}

\markboth{Journal of \LaTeX\ Class Files,~Vol.~14, No.~8, August~2015}
{Shell \MakeLowercase{\textit{et al.}}: Bare Demo of IEEEtran.cls for IEEE Journals}
\maketitle
\begin{abstract}
    We propose a desigining method of a flexible sampling operator for graph signals via a difference-of-convex (DC) optimization algorithm.
    A fundamental challenge in graph signal processing is sampling, especially for graph signals that are not bandlimited.
    In order to sample beyond bandlimited graph signals, there are studies to expand the generalized sampling theory for the graph setting.
    Vertex-wise sampling and flexible sampling are two main strategies to sample graph signals.
    Recovery accuracy of existing vertex-wise sampling methods is highly dependent on specific vertices selected to generate a sampled graph signal that may compromise the accurary especially when noise is generated at the vertices.
    In contrast, a flexible sampling mixes values at multiple vertices to generate a sampled signal for robust sampling; however, existing flexible sampling methods impose strict assumptions and aggressive relaxations.
    To address these limitations, we aim to design a flexible sampling operator without such strict assumptions and aggressive relaxations by introducing DC optimization.
    By formulating the problem of designing a flexible sampling operator as a DC optimization problem, our method ensures robust sampling for graph signals under arbitrary priors based on generalized sampling theory.
    We develop an efficient solver based on the general double-proximal gradient DC algorithm, which guarantees convergence to a critical point. 
    Experimental results demonstrate the superiority of our method in sampling and recovering beyond bandlimited graph signals compared to existing approaches.
\end{abstract}

\begin{IEEEkeywords}
    Graph signal processing, graph signal sampling, generalized sampling theory, difference-of-convex optimization.
\end{IEEEkeywords}

\IEEEpeerreviewmaketitle

\section{Introduction}
%%既存手法と提案手法についてのまとめ
%GSPの概要
\IEEEPARstart{G}{raph} signal processing (GSP) presents an innovative approach to analyzing data depicted as signals situated at the vertices of a graph.
At present, GSP is an active and dynamic domain of investigation within signal processing, encompassing various inquiries ranging from theoretical investigations to real-world implementations.
Recent interest in GSP is to extend classical signal processing theories to the graph setting~\cite{narang2012perfect,agaskar2013spectral,sandryhaila2014big,shuman2015multiscale,tanaka2020gensamp}.
There are several aspects of researches relating to GSP. 
Typical examples include graph learning~\cite{dong2016learning,egilmez2018graph,liu2019graph,dong2019learning}, graph restoration~\cite{ono2015total,onuki2016graph,nagahama2022graph}, and graph signal recovery~\cite{berger2017graph,yamagata2025robust}.
There are also various applications of GSP, including image and point cloud processing~\cite{cheung2018graph,wang2019dynamic,zare2020fusing,takemoto2022graph}, graph neural networks~\cite{bronstein2017geometric}, machine learning~\cite{chen2014semi,gadde2014active}, social networks~\cite{perozzi2014deepwalk}, finantial market~\cite{saboksayr2021learning}, and traffic prediction~\cite{hasanzadeh2019piecewise}.
For an extensive overview of GSP, refer to~\cite{shuman2013emerging,sandryhaila2013discrete,ortega2018graph}.

%GSPにおけるサンプリングの重要性
In the realm of GSP, one of the foundational technologies is \textit{graph signal sampling}~\cite{chen2015discrete, tanaka2020sampling}.
A significant departure from traditional signal sampling is that graph signal sampling has no regular sampling pattern. 
Consequently, extensive research has been devoted to extending the principles of Shannon-Nyquist sampling theory to accommodate graph signals.
Most of the existing methods are based on ensuring bandlimitedness of signals~\cite{chen2015discrete, marques2015sampling, anis2016efficient, tsitsvero2016signals, valsesia2018sampling, tanaka2018spectral, puy2018random, sakiyama2019eigendecomposition, bai2020fast, jayawant2022practical}.
On the other hand, there is a variety of beyond bandlimited graph signals; for instance, meteorological data in mountainous areas, image data of landscapes, and data on water distribution networks~\cite{dong2016learning,cheung2018graph,li2025sensor}.
To sample such graph signals, there are some studies to extend the \textit{generalized sampling theory}~\cite{elder2009beyond, eldar2015sampling} to the graph setting~\cite{chepuri2018graphsamp, tanaka2018spectral, tanaka2020sampling, tanaka2020gensamp,hara2020generalized, hara2021design, hara2022gsss, hara2023graph}.
The theory provides a unified framework for sampling beyond bandlimited signals by assuming priors of signals in the context of the traditional signal sampling.
It aims to achieve not only the perfect recovery but also the best possible recovery based on some strategies such as the least squares or minimax strategies. 

%グラフ信号のサンプリングにおける2つのアプローチ・周波数領域におけるサンプリングの課題
When adapting the generalized sampling theory to graph setting, there are two main approaches for sampling graph signals; sampling in a \textit{graph frequency domain} and sampling in a \textit{graph vertex domain}.
The authors of~\cite{chepuri2018graphsamp} pioneered the extension of the generalized sampling theory to sample and recover beyond bandlimited graph signals, which led to branches of sampling graph signals in the graph frequency domain~\cite{tanaka2020sampling,hara2023graph}. 
Although graph frequency domain methods are elegant, they are computationally expensive due to the need for graph Fourier transforms (GFT) and can be very time-consuming, especially for large graphs, so vertex domain methods have also been explored.

%頂点領域でのサンプリングについての課題
There are some studies about methods sampling graph signals in the graph vertex domain based on the generalized sampling theory~\cite{hara2021design,hara2022gsss}.
Their approaches are to formulate graph signal sampling as an optimization problem to design \textit{sampling operator} that extracts information from an original graph signal to generate a sampled graph signal.
There are two primary strategies for designing a sampling operator: \textit{vertex-wise sampling} and \textit{flexible sampling}.

%vertex-wiseではサンプルされる頂点への依存が高いのでノイズ耐性が低い(のでflexibleなサンプリングが必要)
The method proposed in~\cite{hara2022gsss} adopts a vertex-wise sampling strategy and deals with the arbitrary priors assumed in the generalized sampling theory.
To design a sampling operator, this method employs a greedy algorithm that is computationally efficient and fast, especially for large graphs, as it sequentially selects locally optimal vertices for sampling.
Although the algorithm is efficient, with this strict vertex selection, the recovery accuracy is highly dependent on the vertices selected to generate a sampled signal. 
The recovery accuracy may be compromised especially if noise is generated at the selected vertices, as the effect of the noise may be amplified in the recovery process.

%subspaceで過度な緩和
In contrast, flexible sampling approaches construct a \textit{flexible sampling operator} that mixes signal values across multiple vertices to generate a sampled signal.
A flexible sampling operator reduces dependence on specific vertices and allows for robust sampling and recovery.
It also has the advantage of avoiding the effect of noise at the specific vertices selected to generate a sampled signal, which can be amplified in the recovery process.
To employ this strategy, the authors of~\cite{hara2021design} formulate a problem to design a sampling operator as an optimization problem with an invertibility constraint to achieve the best possible recovery based on the generalized sampling theory.
However, the method introduces a strong assumption of a graph signal and an aggressive relaxation for handling the invertibility constraint to reformulate the original problem into a convex optimization problem.
The assumption leads to the limitation that the method can deal with only part of the priors assumed in the generalized sampling theory, and the relaxation may degrade the recovery accuracy.

%DC最適化によるサンプリングオペレータの設計
Given the limitations of existing sampling methods, particularly those relying on a vertex-wise sampling strategy or convex optimization, a method for designing a flexible sampling operator without applying strong assumptions or aggressive relaxations is desirable.
However, based on the generalized sampling theory, designing a sampling operator that achieves the best possible recovery requires imposing an invertibility constraint.
Since it is difficult to reformulate a problem with this constraint into a convex optimization problem without aggressive relaxations, an approach beyond convex optimization is appropriate to avoid such relaxations.
Here, a question arises:
\textit{How can we formulate a problem with appropriate relaxation of the invertibility constraint to design a flexible sampling operator for sampling beyond bandlimited graph signals under arbitrary priors?}

%提案手法の概要
In this paper, we propose a novel method for designing a flexible sampling operator for sampling graph signals assuming the arbitrary priors in the vertex domain.
We formulate a problem to design a flexible sampling operator with the invertibility constraint to achieve the best possible recovery.
Then, for handling the invertibility constraint, we reformulate the problem as a \textit{difference-of-convex} (DC) optimization problem by introducing a tight relaxation of the constraint.
A DC optimization is a powerful tool for solving optimization problems formulated as the difference of convex functions.
To solve the optimization problem, we develop an efficient solver based on the \textit{general double-proximal gradient DC} (GDPGDC) algorithm~\cite{banert2019general} with guaranteed convergence to a critical point.

The key contributions of our study are summarized below:

\begin{itemize}[align=parleft, left=1.1mm, labelsep=1.1mm, itemindent=0mm]
    \item Formulate the design of a flexible sampling operator that can be applied to arbitrary priors based on the theory as an optimization problem with an invertibility constraint to aim the best possible recovery.
    \item Transform the problem into a tractable DC optimization problem via a tight relaxation of the invertibility constraint using the nuclear norm.
    \item Develop an efficient solver based on the GDPGDC algorithm to solve the DC problem with guaranteed convergence to a critical point.
\end{itemize}

We also conduct sampling and recovering experiments on various types of graph signals to validate the effectiveness of our method.

Table~\ref{table:summary_existing_methods} summarizes the coverage of existing and proposed methods. 
The column “GFT Unrequired” indicates whether the method avoids reliance on a GFT.
The column “Flexible Sampling” specifies if the method adopts a flexible sampling strategy. 
The column “Priors” indicates applicability to the priors assumed in the generalized sampling theory.
\setlength{\textfloatsep}{5pt}
\begin{table}[t]
    \captionsetup{font=small} 
    \centering
    % \caption{\small{Summary of the Coverage of the Existing and Proposed Methods. 
    % The Column "GFT Unrequired" Indicates If the Method Does Not Require the Graph Fourier Transform. 
    % The Column "Flexible Sampling" Indicates Whether the Method Uses Flexible Sampling Strategy. 
    % The Column "Priors" Indicates Whether the Method is Applicable under the Priors Assumed in Generalized Sampling Theory.}}
    \caption{\small{Coverage Summary of Existing and Proposed Methods.}}
    \vspace{-5mm}
    \label{table:summary_existing_methods}
    % \tabcolsep.75\tabcolsep
    \renewcommand{\arraystretch}{1.1}
    \small
    \begin{tabular}[t]{c||c|c|c}
        \toprule
            Method & GFT Unrequired & Flexible Sampling & Priors \\
        \midrule
            \cite{chen2015discrete} & $\checkmark$ & - & BL only \\ \hline
            \cite{tanaka2020sampling} & - & - & Part of \\ \hline
            \cite{jayawant2022practical} & $\checkmark$ & - & BL only \\ \hline
            \cite{hara2021design} & $\checkmark$ & $\checkmark$ & Part of \\ \hline
            % \cite{tanaka2020gensamp} & - & $\checkmark$ & Part of \\ \hline
            \cite{hara2022gsss} & $\checkmark$ & - & $\checkmark$ \\ \hline
            \cite{hara2023graph} & - & $\checkmark$ & $\checkmark$ \\ \hline
            $\textbf{Ours}$ & $\checkmark$ & $\checkmark$ & $\checkmark$ \\
        \bottomrule
    \end{tabular}
\end{table}

This paper is structured as follows.
Section~\ref{sec:preliminaries} reviews the generalized sampling and its extended framework into the graph setting.
The algorithm for solving the DC optimization, GDPGDC, is also introduced in this section.
In Section~\ref{sec:propose}, the formulation of the problem using DC optimization for the proposed method and the algorithm for solving it based on the framework introduced in Section~\ref{sec:preliminaries} are discussed. 
We demonstrate the effectiveness of the proposed method through experiments with existing methods by sampling and recovering graph signals in Section~\ref{sec:experiments}.
Finally, we conclude this paper in Section~\ref{sec:conclusion}.

The preliminary version of this work, which discusses a method of designing a flexible sampling operator by assuming only a limited prior without proposing multiple designs of sampling operators, comprehensive experimental comparisons, and deeper disscussions, has been appeared in a conference proceeding~\cite{yamashita2024apsipa}.

\section{Preliminaries}
\label{sec:preliminaries}
\begin{table}[t]
    \captionsetup{font=small} 
    \centering
    \caption{Notations and Definitions.}
    \vspace{-5mm}
    \label{table:notation}
    \renewcommand{\arraystretch}{1.3}
    \small
    \begin{tabular}[t]{c|c}
        \toprule
        Notations & Definitions \\
        \hline \hline
        $\set{X}$ & a set \\ \hline
        $\vect$ & a vector  \\ \hline
        $\vele, \left[\vect\right]_{ij}$ & the $i$-th element of $\vect$  \\ \hline
        $\Ltwo{\vect}$ & the $\ell_2$ norm of $\vect$, $\Ltwo{\vect} := \sqrt{\sum_i \vele^2}$  \\ \hline
        $\mat$ & a matrix  \\ \hline
        $\ele, \left[\mat\right]_{ij}$ & the $(i,j)$-th element of $\mat$  \\ \hline
        $\tran{\mat}$ & the transpose of $\mat$  \\ \hline
        % $\adj{\mat}$ & adjoint of $\mat$  \\ \hline
        $\inv{\mat}$ & the inverse of $\mat$  \\ \hline
        $\pinv{\mat}$ & the pseudo-inverse of $\mat$  \\ \hline
        $\sval{\mat}{i}$ & the $i$-th singular value of $\mat$  \\ \hline
        $\langle \mat, \matt \rangle$ & \begin{tabular}{c} the inner product of $\mat$ and $\matt$,\\$\langle \mat, \matt \rangle := \sum_{i}\sum_{j} \ele \elee$  \end{tabular}\\ \hline
        $\Lone{\mat}$ & the $\ell_1$ norm of $\mat$, $\Lone{\mat} := \sum_{i}\sum_{j} |\ele|$  \\ \hline
        $\FN{\mat}$ & the Frobenius norm of $\mat$, $\FN{\mat} := \sqrt{\langle \mat, \mat \rangle}$  \\ \hline
        $\NN{\mat}$ & the nuclear norm of $\mat$, $\NN{\mat} := \sum_{i} \sval{\mat}{i}$  \\ \hline
        $\diag{(\cdot)}$ & \begin{tabular}{c}
            a diagonal matrix with $\cdot$\\as its principal diagonal
            \end{tabular}  \\ \hline
        $\indi{\set{C}}{\mat} $ & \begin{tabular}{c}
            the indicator function of a closed set $\set{C}$,\\ $\indi{\set{C}}{\mat} = \begin{cases}
                0 & \text{if } \mat \in \set{C} \\
                \infty & \text{otherwise}
            \end{cases}$ \\
            \end{tabular}  \\ \hline
        $\Expv{\cdot}$ & an expected value \\ \hline
        $\exR$ & an extended real line, $\exR := \realnum \cup \{\infty\}$  \\ \hline
        $\onesN$ & the $N \times N$ identity matrix  \\ \hline
        $\prox_{f}(\cdot)$ & the proximity operator of a convex function $f$  \\
        % $\left[\prox_{f}(\cdot)\right]_{ij}$ & the $(i,j)$-th element of $\prox_{f}(\cdot)$  \\
        \bottomrule
    \end{tabular}
  \end{table}

% Bold lowercase and uppercase letters represent vectors and matrices, respectively. 
% We use $\vele$ and $\ele$ to represent the $i$-th element of a vector $\vect$ and the element in the $i$-th row and $j$-th column of a matrix $\mat$, respectively.
% The $\ell_2$ norm of a vector $\vect$ is denoted by $\Ltwo{\vect} := \sqrt{\sum_i \vele^2}$.
% The transpose and the adjoint of a matrix $\mat$ is denoted as $\tran{\mat}$ and $\adj{\mat}$, respectively. 
% The inverse and the pseudo-inverse of a matrix $\mat$ are denoted by $\inv{\mat}$ and $\pinv{\mat}$, respectively.
% The inner product of two matrices $\mat$ and $\matt$ is denoted by $\langle \mat, \matt \rangle := \sum_{i}\sum_{j} \ele \elee$.
% The $i$-th singular value of $\mat$ is denoted by $\sval{\mat}{i}$.
% The $\ell_1$ norm, Frobenius norm, and the nuclear norm of $\mat$ are denoted by $\Lone{\mat} := \sum_{i}\sum_{j} |\ele|$, $\FN{\mat} := \sqrt{\langle \mat, \mat \rangle}$, and $\NN{\mat} = \sum_i \sval{\mat}{i}$, respectively.
% A diagonal matrix with $\cdot$ as its principal diagonal is denoted by $\diag(\cdot)$.
% The indicator function of a set $\set{C}$ is denoted by $\indi{\set{C}}{\mat}$ that is defined such that $\indi{\set{C}}{\mat} = 0$ when $\mat \in \set{C}$, and $\indi{\set{C}}{\mat} = \infty$ otherwise.
% The extended real line $\exR$ is defined as $\exR := \realnum \cup \{+\infty, -\infty\}$.
% An identity matrix is denoted by $\ones$, and $\onesN$ represents $N \times N$ identity matrix.
This section introduces mathematical tools required for the proposed method, in particular the generalized sampling theory and its extension framework for the graph setting.
The fundamental of GDPDCG, the algorithm needed to solve the proposed optimization problem, is also introduced. 
For references, the notation used in this paper is provided in Table~\ref{table:notation}.

\subsection{Generalized Sampling} \label{sec:gs}
In this subsection, we present the summary of the generalized sampling theory in Hilbert spaces~\cite{eldar2003sampling, eldar2015sampling}.
Let $\sig$ be a vector in a Hilbert space $\HS$ and $\sampsig$ be its samples with the $n$-th sample given by $\elesampsig_n = \inprod{\samp_n}{\sig}$, where $\{ \samp_n \}$ is a Riesz basis and $\inprod{\cdot}{\cdot}$ is an inner product.
The sampled signal $\sampsig$ is given by $\sampsig = \samptranA \sig$, where $\samptranA$ is a sampling operator.
From the sampled signal $\sampsig$, a recovered signal $\recsig$ is obtained by filtering $\sampsig$ with a correction operator $\cortran$ and a reconstruction operator $\rectran$ as follows:
\begin{equation} \label{recsig}
    \recsig = \rectran \cortran \sampsig = \rectran \cortran \samptranA \sig,
\end{equation}
where $\rectran$ corresponds to a basis $\{ w_n \}$ for the reconstruction space, which spans a closed subspace $\set{W}$ of $\HS$.
For achieving a stable recovery, $\{ w_n \}$ satisfies the Riesz basis condition~\cite{eldar2015sampling}.

A reconstruction operator $\rectran$ may be constrained, i.e., it may be predefined for some reasons, such as computational or hardware restrictions.
We call a case where $\rectran$ is constrained as \textit{predefined case} and a case where $\rectran$ is not constrained as an \textit{unconstrained case}.
A correction operator $\cortran$ operates on the samples $\sampsig$ prior before the reconstruction.
The major difference between the generalized sampling theory and the traditional sampling theory for bandlimited signals is the insertion of a correction operator $\cortran$ in the recovery process.
By appropriately designing $\cortran$ with using prior knowledge of the spaces where the original signal $\sig$ belongs to, it is possible to obtain a recovered signal that minimizes the error from the original signal in some senses.

\subsection{The Framework of Graph Signal Sampling}
\label{subsec:gs_sampling}
In this subsection, we provide a brief overview of the generalized sampling theory in the graph setting~\cite{tanaka2020gensamp,tanaka2020sampling}, which forms the basis of our method.
In this paper, we consider a weighted undirected graph $\graph = \{\gver, \gedge\}$, where $\gver$ and $\gedge$ denote a set of vertices and a set of edges between the vertices, respectively. 
The number of vertices is denoted as $N = |\gver|$.
We define an adjacency matrix $\gadjM \in \realNN$, where each element $\eleadjM$ represents the weight of the edge between the $i$-th and $j$-th vertices. 
The weight from $i$th-vertex to $j$th-vertex and the weight from $j$th-vertex to $i$th-vertex are equal for undirected graphs, i.e. $\eleadjM = E_{ji}$.
If the $i$-th and $j$-th vertices are unconnected, then $\eleadjM$ and $E_{ji}$ are set to $0$.
The degree matrix $\gdegM \in \realNN$ is a diagonal matrix, where the $i$-th diagonal element $D_{ii}$ is defined as the sum of weights connected to the vertex $i$, i.e., $D_{ii} = \sum_{j}\eleadjM$.

We use $\gLap := \gdegM - \gadjM$ as a graph Laplacian as a graph variation operator for clarity and specificity.
Since $\gLap$ is a real symmetric matrix, it always admits an eigendecomposition $\gLap = \guniM \geigM \guniMT$, where $\guniM = [\geigvec_1 \dots \geigvec_\numv]$ forms a unitary matrix containing the eigenvectors $\geigvec_{1}, \dots, \geigvec_{\numv}$, and $\geigM = \diag(\geigv_1,\dots,\geigv_\numv)$ comprises the eigenvalues $\geigv_i$.
We denote $\geigv_i$ as the graph frequency, and the graph frequency is smaller as $\geigv_i$ is smaller.

There are two approaches for sampling and recovering: sampling and recovering in a vertex domain~\cite{anis2016efficient,chen2015discrete} and those in a frequency domain~\cite{tanaka2018spectral}.
We describe the sampling approach in a vertex domain as our proposal in this paper focuses on it.

Fig.~\ref{samp_and_rec_framework} shows the outline of sampling and recovering in a vertex domain.
Let $\gsig \in \gset \subseteq \realN$, $\gsampsig \in \realM (\nums \leq \numv)$, and $\grecsig \in \gset$ be an original graph signal, a sampled signal, and a recovered graph signal, respectively.
A graph signal $\gsig$ undergoes sampling by a sampling operator $\sampMT \in \realMN$, i.e., $\gsampsig := \sampMT \gsig$.
Subsequently, the sampled signal $\gsampsig$ is filtered with a correction operator $\corM$ to reduce any errors or distortions introduced during the sampling and recovering process.
Following this, it is further filtered by a reconstruction operator $\recM$ to map the sampled and corrected signal back onto the original graph.
Hence, the recovered signal $\grecsig$ is represented as follows:
\begin{equation}
\label{gsig_recovery}
\grecsig = \recM \corM \gsampsig = \recM \corM \sampMT \gsig.
\end{equation}

The recovery problem entails finding the optimal $\corM$ (and $\recM$ if it is unconstrained) based on assumed priors of the set $\gset$ where $\gsig$ belongs to.
This framework encompasses various situations involving sampling and recovering, including scenarios where the graph signal is bandlimited (refer to~\cite{tanaka2020sampling} for more details).

The correction operator $\corM$ and the reconstruction operator $\recM$ are designed based on well-established strategies; least-squares (LS) strategy, minimax (MX) strategy, and minimum mean squared error (MMSE) strategy.
The LS strategy aims to find the recovered signal $\grecsig$ that minimizes the $\ell_2$ norm of the difference between the sampled signal created from the recovered signal, i.e., $\sampMT \grecsig$ and the sampled signal $\gsampsig$:
\begin{equation} \label{LSstrategy}
\grecsig_{\mathrm{LS}} = \argmin_{\grecsig \in \gset, \ \sampMT \gsig = \gsampsig} \Ltwo{ \sampMT \grecsig - \gsampsig}^{2}.
\end{equation}
The MX strategy attempts to directly control the recovery error $\Ltwo{\grecsig - \gsig}^{2}$ for minimizing the error for the worst feasible signal:
\begin{equation} \label{MXstrategy}
\grecsig_{\mathrm{MX}} = \argmin_{\grecsig \in \gset} \max_{\gsig \in \gset} \Ltwo{ \grecsig - \gsig}^{2}.
\end{equation}
The MMSE strategy aims to minimize the mean squared error (MSE) between $\grecsig$ and $\gsig$:
\begin{equation} \label{MMSEstrategy}
    \grecsig_{\mathrm{MMSE}} = \argmin_{\grecsig \in \gset} \Expv{ \Ltwo{\grecsig - \gsig}^{2}}.
\end{equation}

In graph signal sampling, the sampling operator $\sampMT$ needs to be designed carefully, since the assumptions on $\sampM$ differ between the generalized sampling theory for the classical signal processing and graph signal sampling. 
While the generalized sampling theory typically assumes that $\sampM$ simply selects specific entries from a signal, in graph signal sampling, $\sampM$ must take into account graph structures and reflect dependencies between vertices, leading to a more complex design process for $\sampMT$.
Thus, the sampling operator $\sampMT$ is designed by using $\corM$ and $\recM$ obtained from the strategies described above.

This paper focuses on sampling and recovering graph signals under three representative signal priors assumed in the generalized sampling theory: subspace prior~\cite{tanaka2020sampling}, smoothness prior~\cite{tanaka2020sampling}, and stochastic prior~\cite{hara2023graph}. 
Next, we describe the three priors and provide an overview of previous studies~\cite{tanaka2020gensamp,hara2023graph} related to the design of $\corM$ and $\recM$. 
The designs of the optimal $\corM$ and $\recM$ under these priors are summarized in Table~\ref{summary_corM_recM}. 

%%グラフ信号に関する一般化サンプリングの枠組みの概念図
\begin{figure}[t]
    \captionsetup{font=small}
    \begin{center}
        \begin{minipage}[t]{\hsize}
            \centering
            \includegraphics[width=\hsize]{fig/samp_and_rec.pdf}
        \end{minipage} 
        \vspace{-8mm}
    \end{center}
    \caption{\small{Generalized sampling framework for graph signals, where $\gsig$ and $\grecsig$ are the original and recovered graph signals, respectively.}}
    \label{samp_and_rec_framework}
\end{figure}

% Subspace
\subsubsection{Subspace prior} \label{subsec:pre_sb_priors}
Under the subspace prior, we suppose that a graph signal $\gsig \in \realN$ is characterized by a linear model as follows:
\begin{equation} \label{gs_sb}
    \gsig := \sbgenM \sbexpd,
\end{equation}
where $\sbgenM \in \realNK$ $(\sbgenn \leq \numv)$ is a known generator matrix and $\sbexpd \in \realK$ is an expansion coefficient vector. 
Here, we consider the size of the sampled signal $\nums \leq \sbgenn$ for simplicity.
This formulation encompasses the well-known bandlimited setting.
Specifically, this refers to the case where the generation matrix $\sbgenM$ is a matrix with $\sbgenn$ graph Fourier components corresponding to $\sbgenn$ low graph frequencies, i.e., a matrix constructed by taking the first $\sbgenn$ columns of $\mathbf{U}$.

For the unconstrained case, i.e., $\recM$ is not predefined, the solutions of the LS and MX strategies are same and given by
\begin{equation}\label{sb_un_LS}
    \grecsig = \sbgenM \pinv{(\sampMT \sbgenM)} \sampMT \gsig,
\end{equation}
followed by the correction operator $\corM$ and the reconstruction operator $\recM$ as
\begin{equation}\label{sb_corM_recM}
    \corM = \pinv{(\sampMT \sbgenM)}, \quad \recM = \sbgenM.
\end{equation}

For the predefined case, i.e., $\recM$ is predefined, we take LS and MX strategies for designing the correction operator $\corM$.
The solution of LS strategy is
\begin{equation}\label{sb_sol_LS}
    \grecsig = \recM \pinv{(\sampMT \recM)} \sampMT \gsig,
\end{equation}
where the correction operator $\corM$ is
\begin{equation}\label{sb_corM_LS}
    \corM = \pinv{(\sampMT \recM)},
\end{equation}
and the solution of MX strategy is
\begin{equation}\label{sb_sol_MX}
    \grecsig = \recM (\recMT\recM)^{-1}\recMT\sbgenM \pinv{(\sampMT\sbgenM)} \sampMT \gsig,
\end{equation}
where the correction operator $\corM$ is
\begin{equation}\label{sb_corM_MX}
    \corM = (\recMT\recM)^{-1}\recMT\sbgenM\pinv{(\sampMT\sbgenM)}.
\end{equation}

% Smoothness
\subsubsection{Smoothness Prior} \label{subsec:pre_sm_priors}
The smoothness prior is a less restrictive assumption than the subspace prior that the signal subspace is unknown. 
Specifically, we assume that the signal x is smooth in the following sense: $\Ltwo{\smopeM \gsig}^{2} \leq \smbound^2$ for some constant $\smbound$, where $\smopeM \in \realNN$ is the invertible operator that quantifies the variation of $\gsig$.
Here we consider two approaches to select a solution, the LS and MX strategies, which can be applied in both the unconstrained and predefined cases.

For the unconstrained case, the solution of the LS strategy is given by
\begin{equation}\label{sm_un_sol_LS}
    \grecsig = \smMXrecM \pinv{(\sampMT\smMXrecM)} \sampMT \gsig,
\end{equation}
where $\smMXrecM = (\smopeMT\smopeM)^{-1}\sampM$, with the following $\corM$ and $\recM$:
\begin{equation} \label{sm_un_corM_recM}
    \corM = \pinv{(\sampMT \smMXrecM)}, \quad
    \recM = \smMXrecM = \inv{(\smopeMT\smopeM)}\sampM.
\end{equation}
The solution of the MX strategy coincides with Eq.~\eqref{sm_un_sol_LS}.

For the predefined case, the solution of the LS strategy is given by
\begin{equation}\label{sm_pre_sol_LS}
\grecsig = \smLSrecM \pinv{(\sampMT \smLSrecM)} \sampMT \gsig,
\end{equation}
where $\smLSrecM = \recM \inv{(\recMT \smopeMT \smopeM \recM)} \recMT \sampM$.
Thus, the correction operator $\corM$ is
\begin{equation} \label{sm_pre_LS_corM}
   \corM = (\recMT \smopeMT \smopeM \recM)^{-1} \recMT \sampM \pinv{(\sampMT \smLSrecM)}.
\end{equation}
The solution of the MX strategy is given by
\begin{equation}\label{sm_pre_sol_MX}
    \grecsig = \recM (\recMT\recM)^{-1}\recMT \smMXrecM \pinv{(\sampMT \smMXrecM)} \sampMT \gsig.
\end{equation}
with the following $\corM$:
\begin{equation} \label{sm_pre_MX_corM}
   \corM = (\recMT \recM)^{-1} \recMT \smMXrecM \pinv{(\sampMT \smMXrecM)}.
\end{equation}

\subsubsection{Stochastic Prior} \label{subsec:pre_st_priors}
Under the stochastic prior, we consider that the samples can be obtained with an additive noise $\stnoise \in \realM$ and recover from the noisy sample signal $\ssign = \gsampsig + \stnoise$ as follows: 
\begin{equation} \label{st-reconstruction}
    \grecsig = \recM \corM \ssign = \recM \corM (\gsampsig + \stnoise) = \recM \corM (\sampMT \gsig + \stnoise).
\end{equation}

Suppose a graph signal $\gsig$ as a zero-mean process characterized by graph wide sense stationarity (GWSS)~\cite{hara2023graph}, which parallels the wide sense stationarity (WSS) typically observed in standard signals. 
The graph signal $\gsig$ is associated with a known covariance matrix $\stsigM \in \realNN$. 
Additionally, let $\stnoise$ be another zero-mean GWSS process, possessing a known covariance matrix $\stnoiM \in \realMM$. 
It is important to note that these processes, $\gsig$ and $\stnoise$, are independent of each other. 
Both $\stsigM$ and $\stnoiM$ are autocorrelation matrices in this context.
As for the stochastic prior, we also consider both the unconstrained and predefined cases of the reconstruction operator.

For the unconstrained case, as the reconstruction operator $\recM$ can be freely chosen along with $\corM$. 
The solution of the MMSE strategy is given by
\begin{equation}\label{st_un_MMSE}
    \grecsig = \stsigM \sampM \pinv{(\sampMT \stsigM \sampM + \stnoiM)} (\sampMT \gsig + \stnoise),
\end{equation}
followed by the correction operator $\corM$ and the reconstruction operator $\recM$ as:
\begin{equation} \label{st_un_corM_recM}
    \corM = \pinv{(\sampMT \stsigM \sampM + \stnoiM)}, \quad
    \recM = \stsigM \sampM.
\end{equation}

For the predefined case, the solution of the MMSE strategy is given by
\begin{equation}\label{st_pre_MMSE}
    \grecsig = \recM (\recMT \recM)^{-1} \recMT \stsigM \sampM \pinv{(\sampMT \stsigM \sampM + \stnoiM)} (\sampMT \gsig + \stnoise),
\end{equation}
where the correction operator $\corM$ is
\begin{equation}\label{st_pre_corM}
    \corM = (\recMT \recM)^{-1} \recMT \stsigM \sampM \pinv{(\sampMT \stsigM \sampM + \stnoiM)}.
\end{equation}

%%HとWについてのまとめ
\begin{table*}[t]
    \centering
    % \captionsetup{font=small}
    \captionsetup{font={small, stretch=1.5}}
    % \caption{\small{Designs of the collection operator $\corM$ and the reconstruction operator $\recM$. \\ The following definitions are used: $\smMXrecM = (\smopeMT\smopeM)^{-1}\sampM$ and $\smLSrecM = \recM (\recMT \smopeMT \smopeM \recM)^{-1} \recMT \sampM$}}
    \caption{
        \small
        Designs of the Correction Operator $\corM$ and the Reconstruction Operator $\recM$. \\ 
        \vspace{1mm}
        The Following Definitions are Used: $\smMXrecM = (\smopeMT\smopeM)^{-1}\sampM$ and $\smLSrecM = \recM (\recMT \smopeMT \smopeM \recM)^{-1} \recMT \sampM$.
    }
    \vspace{-5mm}
    \label{summary_corM_recM}
    % \tabcolsep.75\tabcolsep
    \renewcommand{\arraystretch}{1.3}
    \small
    \begin{tabular}[t]{c|c||c|wc{2.5cm}|wc{2.5cm}}
        \toprule
            \multirow{2}{*}{Prior} & \multirow{2}{*}{Criteria} & Predefined & \multicolumn{2}{c}{Unconstrained} \\ \cline{3-5}
            &&$\corM$&$\corM$&$\recM$\\
        \hline
            \multirow{2}{*}{Subspace}
            & LS & $\pinv{(\sampMT \recM)}$ & \multirow{2}{*}{$\pinv{(\sampMT\sbgenM)}$} & \multirow{2}{*}{$\sbgenM$} \\
            & MX & $(\recMT\recM)^{-1}\recMT\sbgenM \pinv{(\sampMT\sbgenM)}$  & & \\
            \hline
            \multirow{2}{*}{Smoothness} 
            & LS & $(\recMT \smopeMT \smopeM \recM)^{-1} \recMT \sampM \pinv{(\sampMT \smLSrecM)}$ & \multirow{2}{*}{$\pinv{(\sampMT \smMXrecM)}$} & \multirow{2}{*}{$\smMXrecM$}\\
            & MX & $(\recMT \recM)^{-1} \recMT \smMXrecM \pinv{(\sampMT \smMXrecM)}$ & & \\
            \hline
            Stochastic & MMSE & $(\recMT\recM)^{-1}\recMT\stsigM\sampM\pinv{(\sampMT\stsigM\sampM + \stnoiM)}$ & $\pinv{(\sampMT\stsigM\sampM + \stnoiM)}$ & $\stsigM\sampM$\\
        \bottomrule
    \end{tabular}
    \vspace{-3mm}
\end{table*}

\subsection{General Double-Proximal Gradient Difference-of-Convex (GDPGDC) Algorithm}
\label{subsec:GDPGDC}
Having discussed the recovery of graph signals based on the generalized sampling theory, we now shift our focus to an algorithm for solving DC minimization problems, which underlie the optimization problems we will formulate.
In this subsection, we introduce the general double-proximal gradient DC (GDPGDC) algorithm~\cite{banert2019general}, a versatile tool for solving DC minimization problems.

GDPGDC algorithm can solve optimization problems in the form of
\begin{equation}
    \label{prob:DC_general}
    \min_{\mat} f_1(\mat) + f_2(\mat) - h(\mattt) \quad \st \ \mattt = \matDC \mat,
\end{equation}
where $f_1: \realnn{n}{m} \rightarrow \realnum$ is a differentiable convex function with a $1/\beta$-Lipschitz continuous gradient for some $\beta > 0$, $f_2:\realnn{n}{m} \rightarrow \exR$ and $h:\realnn{k}{m} \rightarrow \exR$ are proper lower-semicontinuous convex functions, and $\matDC \in \realnn{k}{n}$ is a matrix.

Then, GDPGDC algortihm solves Prob.~\eqref{prob:DC_general} by the following procedures: for $\gamma_1 > 0$ and $\gamma_2 > 0$, iterate
\begin{equation}
    \label{eq:DCAbyGDPG}
    \resizebox{\linewidth-0mm}{!}{$
    \left\lfloor
    \begin{array}{l}
        \mat^{(\Step + 1)} \leftarrow \prox_{\gamma_{1} f_{2}} \left(\mat^{(\Step)} - \gamma_{1} \left(\nabla f_{1} \left(\mat^{(\Step)}\right) - \tran{\matDC} \mathbf{Z}^{(\Step)}\right)\right); \\
        \mattt^{(\Step+1)} \leftarrow \prox_{\gamma_{2}h^{*}} \left(\mattt^{(\Step)} + \gamma_{2}\matDC \mat^{(\Step + 1)}\right);\\
        \Step \leftarrow \Step + 1;
    \end{array}
    \right..
    $}
\end{equation}

Here, the proximity operator of a lower-semicontinuous convex function $f$ with a parameter $\gamma > 0$ is defined as~\cite{bauschke2011convex}
\begin{equation} \label{eq:prox_f2}
\begin{aligned}
    &\prox_{\gamma f}:\realnn{n}{m}\rightarrow\realnn{n}{m}: \\
    &\matt \mapsto \argmin_{\mat} f(\mat) + \frac{1}{2\gamma} \FN{\matt - \mat}^{2}.
\end{aligned}
\end{equation}

The \textit{Fenchel–Rockafellar conjugate function} of $f$ is represented as $f^{*}$ and defined as
\begin{equation}
    f^*(\mat):=\max_{\matt}\langle \mat, \matt \rangle - f(\matt).
\end{equation}
Thanks to Moreau's Identity~\cite[Theorem~14.3(ii)]{bauschke2017correction}, the proximity operator of $f^*$ is calculated with a parameter $\gamma > 0$ as follows:
\begin{equation} \label{eqDCh}
    \prox_{\gamma f^*}(\mat) = \mat - \gamma \prox_{\frac{1}{\gamma} f}\left(\frac{1}{\gamma}\mat\right).
\end{equation}

We summarize the theoretical results for the convergence of GDPGDC algorithm as follows:
\begin{theo}
[{\cite[Proposition~4]{banert2019general}} Convergence of the sequence generated by GDPGDC algorithm]
\label{theo:convergence_GDPGDC}
Let $\inf \{ f_1(\mat) + f_2(\mat) - h(\matDC \mat) \; | \; \mat \in \set{H} \} > - \infty$, and let $0 < \gamma_1 < 2 \beta$ and $0 < \gamma_2 < +\infty$ be satisfied. If $\{\mat^{(\Step)}\}_{\Step \in \mathbb{N}}$ and $\{\mattt^{(\Step)}\}_{\Step \in \mathbb{N}}$, generated by Algorithm~\eqref{eq:DCAbyGDPG}, are bounded, then every cluster point of $\{\mat^{(\Step)}\}_{\Step \in \mathbb{N}}$ is a critical point of Prob.~\eqref{prob:DC_general}.
\end{theo}

\section{Proposed Method}
\label{sec:propose}
The main challenge in graph signal sampling based on the generalized sampling theory is designing the sampling operator $\sampMT$. 
In this section, we present our proposed method for its design.
We formulate the design of a flexible sampling operator as a DC optimization problem.
Our formulation provides a unifying framework applicable to any priors assumed in the generalized sampling theory. 
Finally, we develop an effective solver based on GDPGDC algorithm to solve the problem.

\subsection{The Condition for Designing the Sampling Operator}
\label{subsec:condition_for_sampling_operator}
In this subsection, we discuss about the condition for designing a sampling operator on each prior by using the solutions about recovered siginals introduced in subsection~\ref{subsec:gs_sampling} before discussing the specific designing method for the sampling operator.
In order to achieve the best possible recovery based on Eqs.~\eqref{gsig_recovery} and~\eqref{st-reconstruction}, we try to find $\sampM$ ensuring that the correction operator $\corM$ has the full rank.
To account for the differences in the designs of $\corM$ for each signal prior, we introduce a matrix $\varM$ to find an appropriate $\sampM$ that satisfies the condition that $\corM$ has the full rank.
This matrix $\varM$ is derived based on the conditions of each signal prior.
We will now examine the conditions of $\varM$ on a prior-by-prior basis, and it should be noted that we seek $\sampM$ such that $(\varM \sampM)^\top \varM \sampM$ is invertible as a result.
The summary of the components of the matrix $\varM$ is shown in Table~\ref{table:P}.

\subsubsection{Subspace Prior}
Based on~\eqref{sb_un_LS} and~\eqref{sb_sol_MX}, we seek $\sampM$ for which $\pinv{(\sampMT\sbgenM)}$ has the full column rank, i.e., $\sampMT\sbgenM$ has the full row rank for the unconstrained case and under the MX strategy for the predefined case.
Thus, we need to find $\sampM$ such that $\sampMT \sbgenM (\sampMT \sbgenM)^\top$ is invertible, i.e. $(\sbgenM^\top \sampM)^\top \sbgenMT \sampM$ is invertible.
Consequently, under these scenarios, the matrix $\varM$ results as 
\begin{equation}
    \varM = \sbgenMT.
\end{equation}

Similarly, based on~\eqref{sb_sol_LS}, we seek $\sampM$ for which $\pinv{(\sampMT\recM)}$ has the full rank under the LS strategy for the predefined case.
Thus, we need to find $\sampM$ such that $\sampMT \recM (\sampMT \recM)^\top$ is invertible, i.e. $(\recM^\top \sampM)^\top \recMT \sampM$ is invertible.\footnote{In this paper, for simplicity, we consider the case where $\recM^\top \sampM$ is a wide matrix with $\nums$ rows.}
Consequently, under this senario, the matrix $\varM$ results as 
\begin{equation}
    \varM = \recMT.
\end{equation}

\subsubsection{Smoothness Prior}
Based on~\eqref{sm_un_sol_LS} and~\eqref{sm_pre_sol_MX}, for the unconstrained case and under the MX strategy for the predefined case, we seek $\sampM$ for which $\pinv{(\sampMT \smMXrecM)} = \inv{(\sampMT \smMXrecM)}$, i.e., $\sampMT \smMXrecM$ is invertible.
Let us define the singular value decomposition (SVD) of $\smopeM$ as $\smopeM = \smsvul \smsv \smsvur^\top$.
Then, $\sampMT \smMXrecM$ is transformed as follows:
\begin{align}
    \sampMT \smMXrecM &= \sampMT (\smopeMT\smopeM)^{-1}\sampM  \nonumber \\
    &= \sampMT (\smsvur \smsv^{2} \smsvur^\top)^{-1}\sampM \nonumber \\
    &= (\smsv^{-1} \smsvur^\top \sampM)^\top \smsv^{-1} \smsvur^\top \sampM.
\end{align}
Consequently, one can see that $\sampMT \smMXrecM$ is invertible if and only if $\smsv^{-1} \smsvur^\top \sampM \in \realNM$ has the full column rank so that the matrix $\varM$ results as
\begin{equation}
    \varM = \smsv^{-1} \smsvur^\top.
\end{equation}

Similarly, based on~\eqref{sm_pre_sol_MX}, we seek $\sampM$ for which $\pinv{(\sampMT \smLSrecM)} = \inv{(\sampMT \smLSrecM)}$, i.e., $\sampMT \smLSrecM$ is invertible under the LS strategy for the predefined case.
Let us define the economy-size SVD of $\smopeM \recM \in \realNM$ as $\smopeM \recM = \smsvorul \smsvor \smsvorur^{\top}$, where $\smsvorul \in \realNM$ and $\smsvorur^{\top} \in \realMM$ are the left singular vectors matrix and the transpose of the right singular vectors matrix, respectively, and $\smsvor \in \realMM$ is the singular values matrix, which contains only non-zero singular values.
Then, $\sampMT \smLSrecM$ is transformed as follows:
\begin{align}
    \sampMT \smLSrecM &= \sampMT \recM (\recMT \smopeMT \smopeM \recM)^{-1} \recMT \sampM \nonumber \\
    &= \sampMT \recM \smsvorur \smsvor^{-2} \smsvorur^{\top} \recMT \sampM \nonumber \\
    &= (\smsvor^{-1} \smsvorur^{\top} \recMT \sampM)^{\top} (\smsvor^{-1} \smsvorur^{\top} \recMT \sampM).
\end{align}
Consequently, one can see that $\sampMT \smLSrecM$ is invertible if and only if $\smsvor^{-1} \smsvorur^{\top} \recMT \sampM \in \realNM$ has the full column rank so that the matrix $\varM$ results as
\begin{equation}
    \varM = \smsvor^{-1} \smsvorur^{\top} \recMT.
\end{equation}

\subsubsection{Stochastic Prior}
Based on~\eqref{st_un_MMSE} and~\eqref{st_pre_MMSE}, for both the unconstrained and predefined cases, we seek $\sampM$ for which $\pinv{(\sampMT\stsigM\sampM + \stnoiM)} = \inv{(\sampMT\stsigM\sampM + \stnoiM)}$, i.e., $\sampMT\stsigM\sampM + \stnoiM$ is invertible.
Due to the characteristics of an autocorrelation matrix, $\sampMT\stsigM\sampM + \stnoiM$ is invertible if $\sampMT\stsigM\sampM$ is invertible.
Since $\stsigM$ is symmetric, for some $\stsqM \in \realNN$, $\stsigM$ can be decomposed as
\begin{equation}
    \stsigM = \stsqM^\top \stsqM.
\end{equation}

Then $\sampMT\stsigM\sampM = (\stsqM \sampM)^\top \stsqM \sampM$ is invertible if and only if $\stsqM \sampM \in \realNM$ has the full column rank.
Therefore, the matrix $\varM$ results as 
\begin{equation}
    \varM = \stsqM,
\end{equation}
for both the unconstrained and predefined cases.

%%Pについてのまとめ
\begin{table}[t]
    \captionsetup{font=small}
    \caption{\small{Components of the Matrix $\varM$.}}
    \vspace{-5mm}
    \label{table:P}
    % \tabcolsep.75\tabcolsep
    \renewcommand{\arraystretch}{1.4}
    \small
    \begin{tabular}[t]{c|c||c|c}
        \toprule
            Prior & Criteria & Predefined & Unconstrained\\
            \hline
            \multirow{2}{*}{Subspace} & LS & $\recMT$ & \multirow{2}{*}{$\sbgenM^{\top}$} \\
            & MX & $\sbgenM^{\top}$ & \\
            \hline
            \multirow{2}{*}{Smoothness} & LS & $\smsvor^{-1} \smsvorur^{\top} \recMT$ & \multirow{2}{*}{$\smsv^{-1} \smsvur^{\top}$}\\
            & MX & $\smsv^{-1} \smsvur$ &  \\
            \hline
            Stochastic & MMSE & $\stsqM$ & $\stsqM$ \\
        \bottomrule
    \end{tabular}
\end{table}

\subsection{Problem Formulation}
\label{subsec:prob_formulation}
In this subsection, we discuss about the specific method of desigining a sampling operator $\sampMT$.
Assuming that we have defined $\recM$ and $\corM$ as summarized in Table~\ref{summary_corM_recM}, our current task is to devise an appropriate $\sampM$ that aligns with the recovery process described in~\eqref{gsig_recovery} and~\eqref{st-reconstruction}.

Now, the question arises how to design such $\sampM$. 
To address this, we first formulate a sampling operator design problem as the following problem:
\begin{align}\label{prob:org}
    \min_{\sampM}\fsampM \ \st \
    \begin{cases}
    \sampM \in \sampMSet, \\
    (\varM \sampM)^{\top} \varM \sampM \mbox{ is invertible},
    \end{cases}
\end{align}
where $\fsampM$ is a proper lower-semicontinuous convex function that controls the structure of $\sampM$ such as sparsity or smoothness of its elements, and $\sampMSet$ is a nonempty closed convex set.
The first constraint serves to control the elements of the sampling operator for ensuring stability within the set $\sampMSet$.
The second constraint enforces the full rank in the correction operator $\corM$ for achieving the best possible recovery; however, this constraint is difficult to handle directly.
Fortunately, since the second constraints implies that $\varM \sampM$ has the full rank, we can rewite Prob.~\eqref{prob:org} as follows:
\begin{align}\label{prob:org2}
    \min_{\sampM}\fsampM \ \st \
    \begin{cases}
    \sampM \in \sampMSet, \\
    \varM \sampM \mbox{ has full rank}.
    \end{cases}
\end{align}
Here, since the guarantee of the full rankness of $\varM \sampM$ indicates that the minus of $\rank(\varM \sampM)$ should be as small as possible, we rewrite Prob.~\eqref{prob:org2} as follows:
\begin{align}\label{prob:org3}
    \min_{\sampM} \fsampM - \rank({\varM \sampM}) \ \st \ \sampM\in \sampMSet.
\end{align}
Prob.~\eqref{prob:org3} is still difficult to optimize due to the combinatorial nature of the rank function.
Therefore, by introducing a nuclear norm, which is known for the tightest convex envelope of a rank function~\cite{fazel2002matrix}, we relax the problem as 
\begin{align}\label{prob:org4}
    \min_{\sampM} \fsampM - \NN{\varM \sampM} \ \st \ \sampM\in \sampMSet.
\end{align}

As for the function $\fsampM$ and the set $\sampMSet$, we consider three specific designs. 
\begin{itemize}[align=parleft, left=1.1mm, labelsep=1.1mm, itemindent=0mm]
    \item Design (i): Limit the radius of Frobenius norm of $\sampM$ within $\paramLBall >0$, i.e., 
    \begin{align} \label{prob:Lball}
        \fsampM &= 0, \nonumber \\ 
        \sampMSet &= \LBall := \{ \mat \; | \; \FN{\mat} \leq \paramLBall \}.
    \end{align}
    This allows to control the magnitude of $\sampM$ and prevent the excessive amplification of the graph signal.
    \item Design (ii): Minimize the square of the Frobenius norm of $\sampM$ while its elements are limited between $\lbound$ and $\ubound$, i.e.,
    \begin{align} \label{prob:FNandBox}
        \fsampM &= \paramBox \FN{\sampM}^2, \nonumber \\
        \sampMSet &= \BoxconNM,
    \end{align}
    where $\paramBox$ is a balancing parameter.
    This design encourages a solution where the magnitude is spread out evenly across all elements of $\sampM$ while its elements are constrained to be within the range $[\lbound, \ubound]$.
    In other word, this reduces large values and increases small values that leads to a solution having small fluctuations.
    \item Design (iii): Minimize the $\ell_1$ norm of $\sampM$ while its elements are limited between $\lbound$ and $\ubound$, i.e.,
    \begin{align} \label{prob:L1andBox}
        \fsampM &= \paramBox \Lone{\sampM}, \nonumber \\ 
        \sampMSet &= \BoxconNM,
    \end{align}
    where $\paramBox$ is a balancing parameter.
    This design tends to promote sparsity by minimizing the sum of the absolute values of the elements of $\sampM$, i.e., it encourages $\sampM$ to have many zero or near-zero elements, while its elements are constrained to the range $[\lbound, \ubound]$.
    This approach allows focused sampling to specific vertices, thereby facilitating the effective capture of crucial portions of the graph.
\end{itemize}

\begin{rem}[Trends of sampling operators designed by the proposed method]
    \rm{
        As a flexible sampling operator mixes values on multiple vertices of a graph to generate a sampled signal, it tends to use values on all vertices of the graph to sample a graph signal.
        Thus, there may be a concern that the proposed method is not suitable for applications requiring to sample only a few vertices such as sensor placement problems.
        To address such concerns, a method designing a flexible sampling operator with controlling the number of \textit{sample-contributive vertices} that are used to generate a sampled signal is proposed in~\cite{yamashita2025icassp}.
    }
\end{rem}

\subsection{Optimization}
In this subsection, we describe the optimization algorithm for solving Prob.~\eqref{prob:org4}.
By introducing the indicator function of $\sampMSet$, i.e. $\iota_{\sampMSet}$, which is also a proper lower-semicontinuous convex function, we can reformulate Prob.~\eqref{prob:org4} as follows:
\begin{align}\label{prob:proposed}
    \min_{\sampM} \fsampM + \iota_{\sampMSet}(\sampM) - \NN{\varM \sampM}.
\end{align}

This problem is regarded as the minimization of the difference between two convex functions, and Prob.~\eqref{prob:proposed} is reduced to Prob.~\eqref{prob:DC_general} since $\fsampM + \iota_{\sampMSet}(\sampM)$ and $\NN{\varM \sampM}$ are both proper lower-semicontinuous convex functions.
Thus, GDPGDC algorithm introduced in~\ref{subsec:GDPGDC} can be applied to solve this problem.

The algorithmic procedure for solving the problem is summarized in Algorithm~\ref{algo:GDPG} with the respective functions and variables in~\eqref{prob:DC_general} as follows: $f_1(\mat) = 0$, $f_2(\mat) = \fsamp(\mat) + \indi{\sampMSet}{\mat}$, $h(\mattt) = \NN{\mattt}$, $\mat = \sampM$, and $\matDC = \varM$.

In what follows, we derive specific computations of each step of the algorithm.
For Step~\ref{al_s} of Algorithm~\ref{algo:GDPG}, as there are three specific designs for $\fsamp$ and $\sampMSet$, there are also three patterns of the proximity operator $\prox_{\fsamp + \iota_{\sampMSet}}(\mat)$ as follows:
\begin{itemize}[align=parleft, left=1.1mm, labelsep=1.1mm, itemindent=0mm]
    \item Design (i): Since $\fsampM = 0$ and $\sampMSet = \LBall$ is a nonempty closed convex set, the proximity operator of $\fsamp + \iota_{\sampMSet}$ is equal to the metric projection\footnote{The proximity operator of the indicator function of a nonempty closed convex set $\set{C}$ is identical to the metric projection onto $\set{C}$~\cite{bauschke2011convex}.} onto $\LBall$, i.e.,
    \begin{align}
        \label{eq:proj_Fball}
        \prox_{\fsamp + \iota_{\sampMSet}}(\mat)
        &= \prox_{\iota_{\LBall}} (\mat) \nonumber \\
        &= \begin{cases}
            \mat, & \mathrm{if} \: \mat \in \LBall; \\
            \frac{\paramLBall \mat}{\FN{\mat}}, & \mathrm{otherwise}.
        \end{cases}
    \end{align}
    \item Design (ii): As $\fsampM = \paramBox \FN{\sampM}^2$ and $\sampMSet = \BoxconNM$, the proximity operator $\prox_{\fsamp + \iota_{\sampMSet}}(\mat)$ is calculated as follows: for all $i$ from $1$ to $\numv$ and all $j$ from $1$ to $\nums$,
    \begin{align}
        \label{eq:proj_FNandBox}
        \left[\prox_{\fsamp + \iota_{\sampMSet}}(\mat)\right]_{ij}
        &= \left[\prox_{\gamma( \paramBox \FN{\cdot}^2 + \iota_{\BoxconNM})} (\mat)\right]_{ij} \nonumber \\
        &= \max \left\{ \lbound, \min \left\{ \prox_{\gamma \paramBox \FN{\cdot}^2} (\ele), \ubound \right\}\right\}.
    \end{align}
    For the detail of this caluculation, see Appendix~\ref{appendix:calc_FNandBox}.
    \item Design (iii): As $\fsampM = \paramBox \Lone{\sampM}$ and $\sampMSet = \BoxconNM$, the proximity operator $\prox_{\fsamp + \iota_{\sampMSet}}(\mat)$ is calculated as follows: for all $i$ from $1$ to $\numv$ and all $j$ from $1$ to $\nums$,
    \begin{align}
        \label{eq:proj_L1andBox}
        \left[\prox_{\fsamp + \iota_{\sampMSet}}(\mat)\right]_{ij}
        &= \left[\prox_{\gamma( \paramBox \Lone{\cdot} + \iota_{\BoxconNM})} (\mat)\right]_{ij} \nonumber \\
        &= \max \left\{ \lbound, \min\left\{ \prox_{\gamma \paramBox \Lone{\cdot}}\left(\ele\right), \ubound \right\}\right\}.
    \end{align}
    For the detail of this caluculation, see Appendix~\ref{appendix:calc_L1andBox}.
\end{itemize}

Next, as for Step~\ref{al_z} of Algorithm~\ref{algo:GDPG}, we present the specific computation of the proximity operator of a nuclear norm.
The proximity operator $\prox_{\gamma\NN{\cdot}}(\mat)$ is calculated by
\begin{equation}
  \label{eq:prox_NN1}
  \prox_{\gamma\NN{\cdot}}(\mat)=\mathbf{U}_\mat \mathcal{S}_{\gamma}(\mathbf{\Sigma}_{\mat}) \mathbf{V}_\mat^\top,
\end{equation}
where $\mat = \mathbf{U}_{\mat} \mathbf{\Sigma}_{\mat} \mathbf{V}_{\mat}^\top$ is the SVD of $\mat$, and $\mathcal{S}_{\gamma}(\cdot)$ is the soft-thresholding operator applied to singluar values, defined as:
\begin{equation}
  \mathcal{S}_{\gamma}(\sigma_i) = \max(\sigma_i - \gamma, 0),
\end{equation}
and $\mathcal{S}_{\gamma}(\mathbf{\Sigma}_{\mat})$ is a diagonal matrix with entries $\mathcal{S}_{\gamma}(\sval{\mat}{i})$, i.e., $\mathcal{S}_{\gamma}(\mathbf{\Sigma}_{\mat}) = \diag(\mathcal{S}_{\gamma}(\sval{\mat}{i}))$.
Therefore, the proximity operator of $(\NN{\cdot})^{*}$ in Step~\ref{al_z} in Algorithm~\ref{algo:GDPG} can be calculated with~\eqref{eqDCh} and~\eqref{eq:prox_NN1}.

\begin{algorithm}[t]
    \caption{Sampling operator design algorithm}
    \label{algo:GDPG}
    \begin{algorithmic}[1]
    \renewcommand{\algorithmicrequire}{\textbf{Input:}}
    \renewcommand{\algorithmicensure}{\textbf{Output:}}
    \REQUIRE $ \sampM^{(0)}, \DCw^{(0)}, \paramLBall>0, \gamma_1>0, \gamma_2>0$
    \WHILE{until a stopping criterion is satisfied}
    \STATE $ \sampM^{(\Step+1)} \leftarrow \prox_{\fsamp + \iota_{\sampMSet}}\left( \sampM^{(\Step)} + \gamma_{1} \varM^\top \DCw^{(\Step)}\right)$
    \\by~\eqref{eq:proj_Fball},~\eqref{eq:proj_FNandBox}, or~\eqref{eq:proj_L1andBox}; \label{al_s}
    \STATE $ \DCw^{(\Step+1)} \leftarrow \prox_{\gamma_{2} (\NN{\cdot})^{*}} \left(\DCw^{(\Step)} + \gamma_{2}\varM \sampM^{(\Step + 1)}\right)$
    \\by~\eqref{eqDCh} and~\eqref{eq:prox_NN1}; \label{al_z}
    \STATE $ \Step \leftarrow \Step+1 $;
    \ENDWHILE
    \ENSURE $\sampM^{(\Step)}$
    \end{algorithmic}
\end{algorithm}

\begin{rem}[Convergence of the sequence generated by our algorithm]
\rm{Note that $\prox_{\fsamp + \iota_{\sampMSet}}(\mat)$ always returns a variable whose Frobenius norm is less than $\paramLBall$ or whose elements are between $\lbound$ and $\ubound$.
This indicates that the sequence $\{\sampM^{(\Step)}\}_{\Step \in \mathbb{N}}$ generated by Algorithm~\ref{algo:GDPG} is bounded.
Therefore, from Theorem~\ref{theo:convergence_GDPGDC}, $\{\sampM^{(\Step)}\}_{\Step \in \mathbb{N}}$ is guaranteed to converge to a critical point of Prob.~\eqref{prob:proposed}.}
\end{rem}

\section{Experiments and Results}
\label{sec:experiments}
We demonstrated the effectiveness of our method through sampling and recovering experiments across various types of graph signals.
All experiments were conducted using MATLAB (R2024a) on a Windows 11 computer with Intel Core i9-12900 3.19-GHz processor, 32 GB of RAM, and NVIDIA GeForce RTX 3090.
Our method was compared with the following graph signal sampling methods: NLPD~\cite{chen2015discrete}, SASB~\cite{hara2021design}, AVM~\cite{jayawant2022practical}, GSSS~\cite{hara2022gsss}, and SUST~\cite{hara2023graph}.
NLPD and AVM are representative methods for sampling bandlimited graph signals.
SASB is a method for sampling graph signals under only the subspace prior, which designs a flexible sampling operator.
GSSS is a method for sampling graph signals under arbitrary priors in the graph vertex domain with a vertex-wise strategy.
SUST is a method for sampling graph signals under arbitrary priors in the graph frequency domain.

\subsection{Setup}
We generated random sensor graphs, which are implemented by k nearest neighbor graphs, whose vertices are randomly distributed in 2-D space $[0, 1] \times [0, 1]$, with consisting of $\numv = 256$ vertices by using GSPBox~\cite{perraudin2014gspbox}.
The size of sampled signal was set as $\nums = 16$.
The sampling and recovering framework is illustrated in Fig.~\ref{samp_and_rec_framework}.

We have generated the following types of graph signals.
\begin{itemize}[align=parleft, left=1.1mm, labelsep=1.1mm, itemindent=0mm]
    \item Subspace prior (SB)
    \begin{itemize}[align=parleft, left=0.5mm, labelsep=1.1mm, itemindent=0mm]
        \item Bandlimited (BL) graph signals~\cite{tanaka2020sampling} characterized as:
        \begin{equation} \label{BL}
            \gsig = \sum_{i=1}^{\sbgenn} d_i \geigvec_i = \sbuniBL \sbexpd,
        \end{equation}
        where $\sbuniBL$ is the submatrix of $\guniM$, whose rows are extracted with $\BLset = \{1,\dots,\sbgenn \}$.
        Here, the generator matrix is $\sbgenM = \sbuniBL$;
        \item Periodic graph spectrum (PGS) graph signals with following~\cite{tanaka2020gensamp} that assumes the periodicity of the graph spectrum as follows:
        \begin{equation} \label{PGS}
            \gsig = \guniM \sbgsr \DsampT \sbexpd,
        \end{equation}
        where $\sbgsr$ is a graph spectral response of the generator, which is a diagonal matrix with the $i$-th element $A(\lambda_{i}) = \mathrm{exp}(-1.5 \lambda_{i}/\lambda_{\mathrm{max}})$, where $\lambda_{i}$ and $\lambda_{\mathrm{max}}$ are the $i$-th graph frequency and the largest graph frequency, respectively, and $\Dsamp = [ \onesK \ \onesK \ \dots ] \in \realKN$ is the matrix for the GFT domain upsampling.
        In this case, the generator matrix is $\sbgenM = \guniM \sbgsr \DsampT$;
        \item Piecewise constant (PWC) graph signals~\cite{chen2016representations} that are characterized by constant values in separated vertex regions and are defined as follows with the number of pieces $\sbgenn$:
        \begin{equation} \label{PWC}
            \gsig = \sum_{i=1}^{\sbgenn} d_i \sbpwc{i} = [\sbpwc{1} \ \dots \ \sbpwc{\sbgenn}] \sbexpd,
        \end{equation}
        where $\sbpwc{i} \in \realN$ for any $i = 1,\dots,\sbgenn$ is defined as $\sbelepwc{i}{j} = 1$ when the node $j$ is in the $i$-th piece and $\sbelepwc{i}{j} = 0$ otherwise for any $j = 1,\dots,\numv$.
        In this case, the generator matrix is $\sbgenM = [\sbpwc{1} \ \dots \ \sbpwc{\sbgenn}]$.
    \end{itemize}
    For the signal types under the subspace prior, we set $\sbexpd$ as its elements $d_i \sim \mathcal{N}(1,1)$ for all $i = 1,\dots,\sbgenn$ and $\sbgenn = \nums$.
    \item Smoothness prior (SM)
        \begin{itemize}[align=parleft, left=0.5mm, labelsep=1.1mm, itemindent=0mm]
        \item Gaussian Markov random field (GMRF)~\cite{gadde2015probabilistic} graph signals with the power spectrum for any $i = 1, \ldots, \numv$:
        \begin{equation}
            \smps = 0.1 / \left(\lambda_{i} + 0.1\right);
        \end{equation}
        \item Piesewise linear (PWL) signals that with randomly chosen $8$ vertices having values based on uniformly distributed random numbers in $[-1,1]$ and interpolated values on other vetrices based on the topology of the graph for each intervening vertex.
    \end{itemize}
    \item Stochastic prior (ST)
        \begin{itemize}[align=parleft, left=0.5mm,  labelsep=1.1mm, itemindent=0mm]
        \item GMRF graph signals as stochastic graph signals (SGS) with the power spectrum for any $i = 1, \ldots, \numv$:
        \begin{equation}
            \stelefsigM = \exp \left( - \left(
                \frac{
                    {2 \lambda_{i} - \lambda_{\max}}
                }{
                    {\sqrt{\lambda_{\max}}}
                }
            \right) ^2\right).
        \end{equation}
    \end{itemize}
\end{itemize}

For all graph types above, we also experimented with noisy sampled signals $\ssign := \gsampsig + \stnoise$, where $\stnoise \in \realM$ is generated as a zero-mean white Gaussian noise with its variance $\sigma^2=0.3$.

For the predefined cases, we used the following reconstruction operators for any methods applicable to the predefined cases including our method:
\begin{itemize}[align=parleft, left=1.1mm, labelsep=1.1mm, itemindent=0mm]
    \item SB prior: $\recM = \overline{\sbgenM}$, where each element of $\overline{\sbgenM}$ is obtained by multiplying the corresponding element of $\sbgenM$ by a normal distribution $\mathcal{N}(1,0.1)$.
    \item SM prior: $\recM = (\smopeMT \smopeM)^{-1} \overline{\sampM}$, where each element of $\overline{\sampM}$ is obtained by multiplying the corresponding element of $\sampM$ designed in the unconstrained case by a normal distribution $\mathcal{N}(1,0.1)$.
    \item ST prior: $\recM = \stsigM \overline{\sampM}$, where each element of $\overline{\sampM}$ is obtained by multiplying the corresponding element of $\sampM$ designed in the unconstrained case by a normal distribution $\mathcal{N}(1,0.1)$.
\end{itemize}

As for the parameters, $\paramLBall$ in Prob.~\eqref{prob:Lball} was set as $\paramLBall = \sqrt{\numv\nums}/4$.
The parameters $\lbound$ and $\ubound$ in Prob.~\eqref{prob:FNandBox} and~\eqref{prob:L1andBox} were set as $\lbound = 0$ and $\ubound = 1$.
The parameter $\paramBox$ in Prob.~\eqref{prob:FNandBox} was set as $\paramBox = 0.5$, and $\paramBox$ in Prob.~\eqref{prob:L1andBox} was set as $\paramBox = 0.1$.
For our proposed algorithm, $\gamma_1$ and $\gamma_2$ were set as $\gamma_1 = \gamma_2 = 0.001$ for any optimization designs.
We defined $\sampM^{(0)}$ as a matrix with random Gaussian entries, and the stopping criterion in Argorithm~\ref{algo:GDPG} as $\|\sampM^{(\Step + 1)} - \sampM^{(\Step)}\|_{F}/\|\sampM^{(\Step)}\|_{F} \leq 10^{-5}$.
The parameter settings for the existing methods followed the descriptions in the respective papers.

For sampling graph signals under the smoothness prior, we used the spectral response $\smopefil$ of the smoothness operator $\smopeM$ following~\cite{tanaka2020gensamp} as:
\begin{equation}
    F(\lambda_i) = \frac{\lambda_i}{\lambda_\mathrm{max}} + 1.
\end{equation}

For the quantitative evaluations, we used MSE: $\mathrm{MSE} = \|\tilde{\gsig} - \gsig\|^{2}_2/\numv$, where $\tilde{\gsig}$ and $\gsig$ are the recovered and original graph signals, respectively.
As a note, smaller values of the MSE indicate better recovery accuracy.


\subsection{Results and Discussion}
%% 実験結果のテーブル
\begin{table*}[!t]
    \captionsetup{font=small} 
    \centering
    \caption{Average MSEs in Decibel for 20 Independent Runs.
    UN Refers to the Unconstrained Case, and PRE (LS), PRE (MX), and PRE (MMSE) Refer to the Least-squares, Minimax, and Minimum Mean Square Error Strategy under Predefined Case, Respectively. 
    The Best and Second Best Results in Each Case Are Highlighted in Bold and with Underline, Respectively.}
    \vspace{-2mm}
    \label{table:M16}
    \footnotesize
    \begin{minipage}{1\textwidth}
        \centering
        \begin{tabular}{ccccccccccc}
        \toprule
            Priors & \multirow{2}{*}{\shortstack{Signal\\Types}} & \multirow{2}{*}{Cases} & \multicolumn{7}{c}{Methods} \\
            \cmidrule(lr){4-11}
            &&& NLPD~\cite{chen2015discrete} & SASB~\cite{hara2021design} & AVM~\cite{jayawant2022practical} & GSSS~\cite{hara2022gsss} & SUST~\cite{hara2023graph} & \textbf{Ours (i)} & \textbf{Ours (ii)} & \textbf{Ours (iii)}\\
        \midrule
            \multirow{18}{*}{SB}
                & \multirow{6}{*}{BL}
                    & UN
                    & $-606.576$ & $-315.847$ & $-485.269$ & $\underline{-609.070}$ & $-590.403$ & $\mathbf{-609.860}$ & $-602.645$ & $-584.400$  \\
                    && + noise
                    & $-11.938$ & $10.702$ & $-3.663$ & $-11.212$ & $-34.211$ & $\mathbf{-58.281}$ & $-33.068$ & $\underline{-46.175}$  \\ 
                    && PRE (LS)
                    & $-$ & $-$ & $-$ & $-51.751$ & $-58.977$ & $\mathbf{-59.065}$ & $\underline{-59.043}$ & $-58.804$  \\
                    && + noise
                    & $-$ & $-$ & $-$ & $-11.794$ & $-33.632$ & $\mathbf{-52.618}$ & $-32.648$ & $\underline{-35.837}$  \\ 
                    && PRE (MX)
                    & $-$ & $-$ & $-$ & $-59.066$ & $-59.066$ & $-59.066$ & $-59.066$ & $-59.066$  \\
                    && + noise
                    & $-$ & $-$ & $-$ & $-11.259$ & $-33.797$ & $\mathbf{-52.684}$ & $-32.716$ & $\underline{-44.461}$  \\
                \cmidrule(lr){2-11}
                & \multirow{6}{*}{PGS}
                    & UN
                    & $4.236$ & $-572.889$ & $13.900$ & $\underline{-593.450}$ & $-591.579$ & $\mathbf{-595.307}$ & $-588.880$ & $-588.587$  \\
                    && + noise
                    & $6.094$ & $22.727$ & $11.496$ & $-10.964$ & $\mathbf{-63.596}$ & $-57.510$ & $-42.060$ & $\underline{-62.798}$  \\ 
                    && PRE (LS)
                    & $-$ & $-$ & $-$ & $-32.021$ & $\underline{-41.945}$ & $\mathbf{-42.035}$ & $-41.824$ & $-41.622$  \\
                    && + noise
                    & $-$ & $-$ & $-$ & $-11.397$ & $\mathbf{-41.237}$ & $-40.674$ & $-35.910$ & $\underline{-40.872}$  \\ 
                    && PRE (MX)
                    & $-$ & $-$ & $-$ & $-42.035$ & $-42.035$ & $-42.035$ & $-42.035$ & $-42.035$  \\
                    && + noise
                    & $-$ & $-$ & $-$ & $-10.804$ & $\mathbf{-41.344}$ & $-40.695$ & $-36.066$ & $\underline{-41.281}$  \\ 
                \cmidrule(lr){2-11}
                & \multirow{6}{*}{PWC}
                    & UN
                    & $-4.756$ & $-590.284$ & $6.508$ & $-590.645$ & $-508.225$ & $\mathbf{-593.625}$ & $\underline{-591.624}$ & $-581.720$  \\
                    && + noise
                    & $-2.147$ & $-10.925$ & $13.664$ & $-12.630$ & $69.739$ & $-57.650$ & $\underline{-57.817}$ & $\mathbf{-66.582}$  \\ 
                    && PRE (LS)
                    & $-$ & $-$ & $-$ & $-30.933$ & $-1.229$ & $\mathbf{-39.764}$ & $\underline{-39.745}$ & $-39.673$  \\
                    && + noise
                    & $-$ & $-$ & $-$ & $-12.818$ & $18.139$ & $\underline{-38.697}$ & $-38.622$ & $\mathbf{-39.292}$  \\ 
                    && PRE (MX)
                    & $-$ & $-$ & $-$ & $-39.764$ & $-39.764$ & $-39.764$ & $-39.764$ & $-39.764$  \\
                    && + noise
                    & $-$ & $-$ & $-$ & $-12.333$ & $69.669$ & $-38.730$ & $\underline{-38.749}$ & $\mathbf{-39.380}$  \\  
        \midrule
            \multirow{12}{*}{SM}
                & \multirow{6}{*}{GMRF}
                    & UN
                    & $-14.806$  & $-$ & $5.968$ & $-18.233$ & $-18.907$ & $\mathbf{-21.241}$ & $\underline{-20.696}$ & $-20.150$  \\
                    && + noise
                    & $-6.445$  & $-$ & $11.749$ & $-14.888$ & $-18.804$ & $\mathbf{-21.108}$ & $\underline{-20.653}$ & $-20.117$  \\ 
                    && PRE (LS)
                    & $-$ & $-$ & $-$ & $-18.046$ & $\underline{-18.883}$ & $-16.908$ & $-18.495$ & $\mathbf{-20.139}$  \\
                    && + noise
                    & $-$ & $-$ & $-$ & $-14.484$ & $\underline{-18.774}$ & $-16.861$ & $-18.451$ & $\mathbf{-20.107}$  \\ 
                    && PRE (MX)
                    & $-$ & $-$ & $-$ & $\underline{-18.233}$ & $\mathbf{-18.895}$ & $-18.003$ & $-17.770$ & $-17.659$  \\
                    && + noise
                    & $-$ & $-$ & $-$ & $-14.888$ & $\mathbf{-18.792}$ & $\underline{-18.012}$ & $-17.652$ & $-17.642$  \\ 
            \cmidrule(lr){2-11}
            & \multirow{6}{*}{PWL}
                    & UN
                    & $\underline{-62.276}$  & $-$ & $-37.465$ & $-34.056$ & $-35.545$ & $\mathbf{-67.558}$ & $-53.508$ & $-44.042$  \\
                    && + noise
                    & $-11.772$  & $-$ & $14.800$ & $-22.281$ & $-34.814$ & $\mathbf{-55.439}$ & $\underline{-51.715}$ & $-43.563$  \\ 
                    && PRE (LS)
                    & $-$ & $-$ & $-$ & $-30.232$ & $\mathbf{-30.314}$ & $-30.138$ & $-30.182$ & $\underline{-30.275}$  \\
                    && + noise
                    & $-$ & $-$ & $-$ & $-29.800$ & $\mathbf{-30.302}$ & $-30.140$ & $-30.180$ & $\underline{-30.273}$  \\
                    && PRE (MX)
                    & $-$ & $-$ & $-$ & $-30.552$ & $-31.089$ & $\mathbf{-35.149}$ & $\underline{-34.955}$ & $-32.990$  \\
                    && + noise
                    & $-$ & $-$ & $-$ & $-26.006$ & $-30.901$ & $\mathbf{-33.783}$ & $\underline{-33.552}$ & $-32.653$  \\ 
            \midrule
            \multirow{4}{*}{ST}
                & \multirow{4}{*}{SGS}
                    & UN
                    & $-2.998$  & $-$ & $6.618$ & $-8.582$ & $-8.918$ & $\mathbf{-9.411}$ & $\underline{-9.389}$ & $-9.097$  \\
                    && + noise
                    & $-0.386$  & $-$ & $15.956$ & $-8.017$ & $-8.808$ & $\mathbf{-9.385}$ & $-8.887$ & $\underline{-9.070}$  \\ 
                    && PRE (MMSE)
                    & $-$ & $-$ & $-$ & $-8.579$ & $-8.903$ & $\mathbf{-9.405}$ & $\underline{-9.383}$ & $-9.086$  \\
                    && + noise
                    & $-$ & $-$ & $-$ & $-8.013$ & $-8.792$ & $\mathbf{-9.378}$ & $-8.885$ & $\underline{-9.060}$  \\ 
        \bottomrule
        \end{tabular}
        \\ \vspace*{1mm} 
        \small
        \raggedright
    \end{minipage}
    \vspace*{-4mm}
\end{table*}
% the result of experiments
\def\resultA{-52.324}
\def\resultB{-28.966}
\def\resultC{-17.603}
\def\resultD{-35.036}
\def\resultE{-33.830}
\def\resultF{-56.067}
\def\resultG{-53.791}
\def\resultH{-46.283}
\def\resultI{-29.306}
\def\resultJ{-29.215}
\def\resultK{-29.095}
\def\resultL{-29.198}
\def\resultM{-29.265}
\def\resultN{-29.630}
\def\resultO{-29.901}
\def\resultP{-36.913}
\def\resultQ{-34.960}
\def\resultR{-32.282}

% the result of experiments (with noise)
\def\resultNoiseA{-10.219}
\def\resultNoiseB{-28.966}
\def\resultNoiseC{30.637}
\def\resultNoiseD{-20.045}
\def\resultNoiseE{-32.793}
\def\resultNoiseF{-49.925}
\def\resultNoiseG{-51.806}
\def\resultNoiseH{-45.605}
\def\resultNoiseI{-28.734}
\def\resultNoiseJ{-29.176}
\def\resultNoiseK{-29.106}
\def\resultNoiseL{-29.197}
\def\resultNoiseM{-29.264}
\def\resultNoiseN{-24.179}
\def\resultNoiseO{-29.712}
\def\resultNoiseP{-32.720}
\def\resultNoiseQ{-32.303}
\def\resultNoiseR{-31.777}

\newcounter{subsubfigure}

\newcommand{\subsubcaption}[2]{%
    \backsubcaption{#1}
    \stepcounter{subsubfigure}%
    \subcaption[\alph{subfigure}-\arabic{subsubfigure}]{#2}%
    \resetsubsubcaption
}

\newcommand{\backsubcaption}[1]{%
    \addtocounter{subfigure}{-#1}%
}

\newcommand{\resetsubsubcaption}{%
    \setcounter{subsubfigure}{1}%
}

\begin{figure*}[t]   
    \centering
    \captionsetup{font=small} 

    % 行1
    \begin{subfigure}[t]{.115\textwidth}
        \includegraphics[width=\linewidth]{fig/PWL/01withoutNoise/00original.pdf}
        \vspace*{-6mm}
        % \resetsubsubcaption
        \subcaption{Original \\ MSE [dB] \\ (Case)}
    \end{subfigure}
    \renewcommand{\thesubfigure}{\alph{subfigure}-\arabic{subsubfigure}}
    \hfill
    \begin{subfigure}[t]{.115\textwidth}
        \includegraphics[width=\linewidth]{fig/PWL/01withoutNoise/01NLPD.pdf}
        \vspace*{-6mm}
        \resetsubsubcaption
        \subcaption{NLPD~\cite{chen2015discrete} \\ \resultA \\ (UN) }
    \end{subfigure}
    \hfill
    \begin{subfigure}[t]{.115\textwidth}
        \includegraphics[width=\linewidth]{fig/PWL/01withoutNoise/03AVM.pdf}
        \vspace*{-6mm}
        \subcaption{AVM~\cite{jayawant2022practical} \\ \resultC \\ (UN) }
    \end{subfigure}
    \hfill
    \begin{subfigure}[t]{.115\textwidth}
        \includegraphics[width=\linewidth]{fig/PWL/01withoutNoise/04GSSS-UN.pdf}
        \vspace*{-6mm}
        \subcaption{GSSS~\cite{hara2022gsss} \\ \resultD \\ (UN) }
    \end{subfigure}
    \hfill
    \begin{subfigure}[t]{.115\textwidth}
        \includegraphics[width=\linewidth]{fig/PWL/01withoutNoise/05SUST-UN.pdf}
        \vspace*{-6mm}
        \subcaption{SUST~\cite{hara2023graph} \\ \resultE \\ (UN)}
    \end{subfigure}
    \hfill
    \begin{subfigure}[t]{.115\textwidth}
        \includegraphics[width=\linewidth]{fig/PWL/01withoutNoise/06Design1-UN.pdf}
        \vspace*{-6mm}
        \subcaption{\textbf{Ours (i)} \\ \textbf{\resultF} \\ (UN)}
    \end{subfigure}
    \hfill
    \begin{subfigure}[t]{.115\textwidth}
        \includegraphics[width=\linewidth]{fig/PWL/01withoutNoise/07Design2-UN.pdf}
        \vspace*{-6mm}
        \subcaption{\textbf{Ours (ii)} \\ \underline{\resultG} \\ (UN)}
    \end{subfigure}
    \hfill
    \begin{subfigure}[t]{.115\textwidth}
        \includegraphics[width=\linewidth]{fig/PWL/01withoutNoise/08Design3-UN.pdf}
        \vspace*{-6mm}
        \subcaption{\textbf{Ours (iii)} \\ \resultH \\ (UN)}
    \end{subfigure}
    \hfill
        \begin{subfigure}{.024\textwidth}
        \includegraphics[width=\linewidth]{fig/PWL/01withoutNoise/colorbar.pdf}
        % 最後の画像にはサブキャプションをつけない
    \end{subfigure}

    % 行2
    \begin{subfigure}[t]{.115\textwidth}
        \hspace*{\fill}
    \end{subfigure}
    \hfill
    \begin{subfigure}[t]{.115\textwidth}
        \includegraphics[width=\linewidth]{fig/PWL/02withNoise/01NLPD_noise.pdf}
        \vspace*{-6mm}
        \subsubcaption{7}{NLPD~\cite{chen2015discrete} \\ \resultNoiseA \\ (UN+noise) }
    \end{subfigure}
    \hfill
    \begin{subfigure}[t]{.115\textwidth}
        \includegraphics[width=\linewidth]{fig/PWL/02withNoise/03AVM_noise.pdf}
        \vspace*{-6mm}
        \subsubcaption{0}{AVM~\cite{jayawant2022practical} \\ \resultNoiseC \\ (UN+noise) }
    \end{subfigure}
    \hfill
    \begin{subfigure}[t]{.115\textwidth}
        \includegraphics[width=\linewidth]{fig/PWL/02withNoise/04GSSS-UN_noise.pdf}
        \vspace*{-6mm}
        \subsubcaption{0}{GSSS~\cite{hara2022gsss} \\ \resultNoiseD \\ (UN+noise) }
    \end{subfigure}
    \hfill
    \begin{subfigure}[t]{.115\textwidth}
        \includegraphics[width=\linewidth]{fig/PWL/02withNoise/05SUST-UN_noise.pdf}
        \vspace*{-6mm}
        \subsubcaption{0}{SUST~\cite{hara2023graph} \\ \resultNoiseE \\ (UN+noise)}
    \end{subfigure}
    \hfill
    \begin{subfigure}[t]{.115\textwidth}
        \includegraphics[width=\linewidth]{fig/PWL/02withNoise/06Design1-UN_noise.pdf}
        \vspace*{-6mm}
        \subsubcaption{0}{\textbf{Ours (i)} \\ \underline{\resultNoiseF} \\ (UN+noise)}
    \end{subfigure}
    \hfill
    \begin{subfigure}[t]{.115\textwidth}
        \includegraphics[width=\linewidth]{fig/PWL/02withNoise/07Design2-UN_noise.pdf}
        \vspace*{-6mm}
        \subsubcaption{0}{\textbf{Ours (ii)} \\ \textbf{\resultNoiseG} \\ (UN+noise)}
    \end{subfigure}
    \hfill
    \begin{subfigure}[t]{.115\textwidth}
        \includegraphics[width=\linewidth]{fig/PWL/02withNoise/08Design3-UN_noise.pdf}
        \vspace*{-6mm}
        \subsubcaption{0}{\textbf{Ours (iii)} \\ \resultNoiseH \\ (UN+noise)}
    \end{subfigure}
    \hfill
        \begin{subfigure}{.024\textwidth}
        \includegraphics[width=\linewidth]{fig/PWL/01withoutNoise/colorbar.pdf}
        % 最後の画像にはサブキャプションをつけない
    \end{subfigure}

    \vspace*{-1mm}
    \caption{An example of PWL graph signals under the smoothness prior defined on a sensor graph and its sampled and recovered signals using each method under the unconstrained case with $\numv=256, \nums=16$.
    The color of each vertex indicates the magnitude of the signal value.
    UN refers the recovered signals in the unconstrained case.
    Ours (i)-(iii) refer to the Design (i)-(iii) of the proposed method described in the section~\ref{sec:propose}, respectively.
    The best and second best results in each case are highlighted in bold and with underline, respectively.}
    % \\*Note: Smaller values indicate better recovery accuracy.
\label{fig:results_PWL_UN}
\vspace*{-4mm}
\end{figure*}
% the result of experiments
\def\resultA{4.044}
\def\resultB{-573.062}
\def\resultC{14.327}
\def\resultD{-586.704}
\def\resultE{-577.459}
\def\resultF{-593.891}
\def\resultG{-595.600}
\def\resultH{-569.083}
\def\resultI{-25.582}
\def\resultJ{-23.369}
\def\resultK{-35.543}
\def\resultL{-35.523}
\def\resultM{-35.85}
\def\resultN{-35.543}
\def\resultO{-35.543}
\def\resultP{-35.543}
\def\resultQ{-35.543}
\def\resultR{-35.543}

% the result of experiments (with noise)
\def\resultNoiseA{3.270}
\def\resultNoiseB{-13.828}
\def\resultNoiseC{21.271}
\def\resultNoiseD{-17.747}
\def\resultNoiseE{-40.843}
\def\resultNoiseF{-62.137}
\def\resultNoiseG{-60.711}
\def\resultNoiseH{-66.376}
\def\resultNoiseI{-10.155}
\def\resultNoiseJ{-22.455}
\def\resultNoiseK{-35.159}
\def\resultNoiseL{-34.975}
\def\resultNoiseM{-35.183}
\def\resultNoiseN{-16.762}
\def\resultNoiseO{-31.797}
\def\resultNoiseP{-35.149}
\def\resultNoiseQ{-35.081}
\def\resultNoiseR{-35.299}

\newcommand{\addsubsubcaption}{%
    \stepcounter{subsubfigure}
}

\begin{figure*}[!ht]   
    \centering
    \captionsetup{font=small} 

    % 行1
    \begin{subfigure}[t]{.025\textwidth}
        \hspace*{\fill}
    \end{subfigure}
    \begin{subfigure}[t]{.139\textwidth}
        \includegraphics[width=\linewidth]{fig/PWC/01withoutNoise/00original.pdf}
        \vspace*{-6mm}
        % \resetsubsubcaption
        \subcaption{Original \\ MSE [dB] \\ (Case)}
    \end{subfigure}
    \renewcommand{\thesubfigure}{\alph{subfigure}-\arabic{subsubfigure}}
    \hfill
    \begin{subfigure}[t]{.139\textwidth}
        \includegraphics[width=\linewidth]{fig/PWC/01withoutNoise/09GSSS-LS.pdf}
        \vspace*{-6mm}
        \resetsubsubcaption
        \subcaption{GSSS~\cite{hara2022gsss} \\ \resultI \\ (LS) }
    \end{subfigure}
    \hfill
    \begin{subfigure}[t]{.139\textwidth}
        \includegraphics[width=\linewidth]{fig/PWC/01withoutNoise/10SUST-LS.pdf}
        \vspace*{-6mm}
        \subcaption{SUST~\cite{hara2023graph} \\ \resultJ \\ (LS)}
    \end{subfigure}
    \hfill
    \begin{subfigure}[t]{.139\textwidth}
        \includegraphics[width=\linewidth]{fig/PWC/01withoutNoise/11Design1-LS.pdf}
        \vspace*{-6mm}
        \subcaption{\textbf{Ours (i)} \\ \textbf{\resultK} \\ (LS)}
    \end{subfigure}
    \hfill
    \begin{subfigure}[t]{.139\textwidth}
        \includegraphics[width=\linewidth]{fig/PWC/01withoutNoise/12Design2-LS.pdf}
        \vspace*{-6mm}
        \subcaption{\textbf{Ours (ii)} \\ \underline{\resultL} \\ (LS)}
    \end{subfigure}
    \hfill
    \begin{subfigure}[t]{.139\textwidth}
        \includegraphics[width=\linewidth]{fig/PWC/01withoutNoise/13Design3-LS.pdf}
        \vspace*{-6mm}
        \subcaption{\textbf{Ours (iii)} \\ \resultM \\ (LS)}
    \end{subfigure}
    \hfill
        \begin{subfigure}{.0215\textwidth}
        \includegraphics[width=\linewidth]{fig/PWC/01withoutNoise/colorbar.pdf}
        % 最後の画像にはサブキャプションをつけない
    \end{subfigure}
    \begin{subfigure}[t]{.025\textwidth}
        \hspace*{\fill}
    \end{subfigure}

    % 行2
    \begin{subfigure}[t]{.025\textwidth}
        \hspace*{\fill}
    \end{subfigure}
    \begin{subfigure}[t]{.139\textwidth}
        \hspace*{\fill}
    \end{subfigure}
    \hfill
    \begin{subfigure}[t]{.139\textwidth}
        \includegraphics[width=\linewidth]{fig/PWC/02withNoise/09GSSS-LS_noise.pdf}
        \vspace*{-6mm}
        \subsubcaption{5}{GSSS~\cite{hara2022gsss} \\ \resultNoiseI \\ (LS+noise)}
    \end{subfigure}
    \hfill
    \begin{subfigure}[t]{.139\textwidth}
        \includegraphics[width=\linewidth]{fig/PWC/02withNoise/10SUST-LS_noise.pdf}
        \vspace*{-6mm}
        \subsubcaption{0}{SUST~\cite{hara2023graph} \\ \resultNoiseJ \\ (LS+noise)}
    \end{subfigure}
    \hfill
    \begin{subfigure}[t]{.139\textwidth}
        \includegraphics[width=\linewidth]{fig/PWC/02withNoise/11Design1-LS_noise.pdf}
        \vspace*{-6mm}
        \subsubcaption{0}{\textbf{Ours (i)}\\\underline{\resultNoiseK} \\ (LS+noise)}
    \end{subfigure}
    \hfill
    \begin{subfigure}[t]{.139\textwidth}
        \includegraphics[width=\linewidth]{fig/PWC/02withNoise/12Design2-LS_noise.pdf}
        \vspace*{-6mm}
        \subsubcaption{0}{\textbf{Ours (ii)}\\\resultNoiseL \\ (LS+noise)}
    \end{subfigure}
    \hfill
    \begin{subfigure}[t]{.139\textwidth}
        \includegraphics[width=\linewidth]{fig/PWC/02withNoise/13Design3-LS_noise.pdf}
        \vspace*{-6mm}
        \subsubcaption{0}{\textbf{Ours (iii)}\\\textbf{\resultNoiseM} \\ (LS+noise)}
    \end{subfigure}
    \hfill
        \begin{subfigure}{.0215\textwidth}
        \includegraphics[width=\linewidth]{fig/PWC/01withoutNoise/colorbar.pdf}
        % 最後の画像にはサブキャプションをつけない
    \end{subfigure}
    \begin{subfigure}[t]{.025\textwidth}
        \hspace*{\fill}
    \end{subfigure}
    
    % 行3
    \begin{subfigure}[t]{.025\textwidth}
        \hspace*{\fill}
    \end{subfigure}
    \begin{subfigure}[t]{.139\textwidth}
        \hspace*{\fill}
    \end{subfigure}
    \hfill
    \begin{subfigure}[t]{.139\textwidth}
        \includegraphics[width=\linewidth]{fig/PWC/01withoutNoise/14GSSS-MX.pdf}
        \vspace*{-6mm}
        \addsubsubcaption
        \subsubcaption{5}{GSSS~\cite{hara2022gsss} \\ \resultN \\ (MX)}
    \end{subfigure}
    \hfill
    \begin{subfigure}[t]{.139\textwidth}
        \includegraphics[width=\linewidth]{fig/PWC/01withoutNoise/15SUST-MX.pdf}
        \vspace*{-6mm}
        \addsubsubcaption
        \subsubcaption{0}{SUST~\cite{hara2023graph} \\ \resultO \\ (MX)}
    \end{subfigure}
    \hfill
    \begin{subfigure}[t]{.139\textwidth}
        \includegraphics[width=\linewidth]{fig/PWC/01withoutNoise/16Design1-MX.pdf}
        \vspace*{-6mm}
        \addsubsubcaption
        \subsubcaption{0}{\textbf{Ours (i)} \\ \resultP \\ (MX)}
    \end{subfigure}
    \hfill
    \begin{subfigure}[t]{.139\textwidth}
        \includegraphics[width=\linewidth]{fig/PWC/01withoutNoise/17Design2-MX.pdf}
        \vspace*{-6mm}
        \addsubsubcaption
        \subsubcaption{0}{\textbf{Ours (ii)} \\ \resultQ \\ (MX)}
    \end{subfigure}
    \hfill
    \begin{subfigure}[t]{.139\textwidth}
        \includegraphics[width=\linewidth]{fig/PWC/01withoutNoise/18Design3-MX.pdf}
        \vspace*{-6mm}
        \addsubsubcaption
        \subsubcaption{0}{\textbf{Ours (iii)} \\ \resultR \\ (MX)}
    \end{subfigure}
    \hfill
        \begin{subfigure}{.0215\textwidth}
        \includegraphics[width=\linewidth]{fig/PWC/01withoutNoise/colorbar.pdf}
        % 最後の画像にはサブキャプションをつけない
    \end{subfigure}
    \begin{subfigure}[t]{.025\textwidth}
        \hspace*{\fill}
    \end{subfigure}

    % 行4
    \begin{subfigure}[t]{.025\textwidth}
        \hspace*{\fill}
    \end{subfigure}
    \begin{subfigure}[t]{.139\textwidth}
        \hspace*{\fill}
    \end{subfigure}
    \hfill
    \begin{subfigure}[t]{.139\textwidth}
        \includegraphics[width=\linewidth]{fig/PWC/02withNoise/14GSSS-MX_noise.pdf}
        \vspace*{-6mm}
        \addsubsubcaption \addsubsubcaption
        \subsubcaption{5}{GSSS~\cite{hara2022gsss} \\ \resultNoiseN \\ (MX+noise)}
    \end{subfigure}
    \hfill
    \begin{subfigure}[t]{.139\textwidth}
        \includegraphics[width=\linewidth]{fig/PWC/02withNoise/15SUST-MX_noise.pdf}
        \vspace*{-6mm}
        \addsubsubcaption \addsubsubcaption
        \subsubcaption{0}{SUST~\cite{hara2023graph} \\ \resultNoiseO \\ (MX+noise)}
    \end{subfigure}
    \hfill
    \begin{subfigure}[t]{.139\textwidth}
        \includegraphics[width=\linewidth]{fig/PWC/02withNoise/16Design1-MX_noise.pdf}
        \vspace*{-6mm}
        \addsubsubcaption \addsubsubcaption
        \subsubcaption{0}{\textbf{Ours (i)} \\ \underline{\resultNoiseP} \\ (MX+noise)}
    \end{subfigure}
    \hfill
    \begin{subfigure}[t]{.139\textwidth}
        \includegraphics[width=\linewidth]{fig/PWC/02withNoise/17Design2-MX_noise.pdf}
        \vspace*{-6mm}
        \addsubsubcaption \addsubsubcaption
        \subsubcaption{0}{\textbf{Ours (ii)} \\ \resultNoiseQ \\ (MX+noise)}
    \end{subfigure}
    \hfill
    \begin{subfigure}[t]{.139\textwidth}
        \includegraphics[width=\linewidth]{fig/PWC/02withNoise/18Design3-MX_noise.pdf}
        \vspace*{-6mm}
        \addsubsubcaption \addsubsubcaption
        \subsubcaption{0}{\textbf{Ours (iii)} \\ \textbf{\resultNoiseR} \\ (MX+noise)}
    \end{subfigure}
    \hfill
    \begin{subfigure}{.0215\textwidth}
        \includegraphics[width=\linewidth]{fig/PWC/01withoutNoise/colorbar.pdf}
        % 最後の画像にはサブキャプションをつけない
    \end{subfigure}
    \begin{subfigure}[t]{.025\textwidth}
        \hspace*{\fill}
    \end{subfigure}

    \vspace*{-1mm}
    \caption{An example of PWC graph signals under the subspace prior defined on a sensor graph and its sampled and recovered signals using each method under the predefined case with $\numv=256, \nums=16$.
    The color of each vertex indicates the magnitude of the signal value.
    LS and MX refer the recovered signals under the least-squares and minimax strategy in the predefined case, respectively.
    Ours (i)-(iii) refer to the Design (i)-(iii) of the proposed method described in the section~\ref{sec:propose}, respectively.
    The best and second best results in each case are highlighted in bold and with underline, respectively.}
    % \\*Note: Smaller values indicate better recovery accuracy.
\label{fig:results_PWC_PRE}
\vspace*{-5mm}
\end{figure*}
\begin{figure}[t]   
    \centering
    \captionsetup{font=small} 
    \begin{subfigure}[t]{0.312\linewidth}
        \centering
        \includegraphics[width=\linewidth]{fig/sampM/ST/06design1.pdf}
        \vspace*{-6mm}
        \subcaption{Design (i)}
    \end{subfigure}
    \hfill
    \begin{subfigure}[t]{0.32\linewidth}
        \centering
        \includegraphics[width=\linewidth]{fig/sampM/ST/07design2_log.pdf}
        \vspace*{-6mm}
        \subcaption{Design (ii)}
    \end{subfigure}
    \hfill
    \begin{subfigure}[t]{0.32\linewidth}
        \centering
        \includegraphics[width=\linewidth]{fig/sampM/ST/08design3_log.pdf}
        \vspace*{-6mm}
        \subcaption{Design (iii)}
    \end{subfigure}
    \vspace*{-1mm}
    \caption{An example of visualized sampling operators designed by the proposed method for a graph signal generated under the stochastic prior in the experiment under the unconstrained case. The color of each cell indicates the value of each element of the sampling operators.}
    \label{fig:results_sampM}
    \vspace*{-1mm}
\end{figure}

Table~\ref{table:M16} represents the averaged MSEs between original graph signals and recovered graph signals generated by each method in decibels obtained from 20 independent runs for each graph signal type.
Fig.~\ref{fig:results_PWL_UN} and~\ref{fig:results_PWC_PRE} visualize examples of graph signals and its sampled and recovered signals by each method. 
Fig.~\ref{fig:results_PWL_UN} visualizes examples of the unconstrained case, and Fig.~\ref{fig:results_PWC_PRE} visualizes examples of the predefined case.

Ours (i)-(iii) refer to the Design (i)-(iii) of the proposed method described in the section~\ref{sec:propose}, respectively.
Since NLPD, SASB, and AVM are not designed separately to the unconstrained or predefined case, they are indicated only in the UN rows in the table and not shown in Fig.~\ref{fig:results_PWC_PRE}.
SASB is also not applicable to the smoothness and stochastic priors since it is designed to sample only graph signals under the subspace prior, so those rows in the table are left blank.
The recovery results in the table and the figures are expressed in decibels that lower numbers indicating better results.
The best and second best results in each case are highlighted in bold and with underline, respectively.

From the table, the results of our method are at least in the two best results for most of the cases regardless of the graph types and even with or without the noise on the sampled signals.
Especially, for the unconstrained cases in any graph types under any priors, either design of our method outperforms the existing methods.

For the predefined cases, although the results in MSE are still competitive, there are cases that the results of our method are not the best and those of the SUST are the best.
This is because SUST is designed to sample in the frequency domain so that it also mixes the values on multiple vertices to generate a sampled signal.
Thus, the results of SUST are relatively stable and competitive for most types of signals regardless of the sampled signals with or without noise.
However, because SUST uses a low-pass filter and shows its strength when sampling a smooth signal, it is heavily affected by noise and produces unfavourable results when sampling signals with abrupt changes in values, such as PWC graph signals.

As for the PWL graph signal, NLPD shows relatively good results in MSE, especially in the noiseless cases, which can also be seen in (b-1) in Fig.~\ref{fig:results_PWL_UN}.
This is because that PWL graph signals have similar behaviors to BL graph signals in the frequency domain.
However, when the noise is added to the sampled signals, the results of our method are superior to that of NLPD, which can be seen in the bottom row of Fig.~\ref{fig:results_PWL_UN}.
 
For any signal types under the subspace prior, the results of GSSS, SUST, and our method in MSE are identical under the predefined cases with the MX strategy, which can also be seen visually in (b-3) - (f-3) in Fig.~\ref{fig:results_PWC_PRE}.
This is because all these methods aim to design $\sampM$ such that $\sampMT \sbgenM$ is invertible under the setting of this experiment, and as a result, as can be seen from Eqs.~\eqref{gs_sb},~\eqref{sb_sol_MX}, and~\eqref{sb_corM_MX}, the effect of $\sampM$ on the recovered signal is eliminated.
Though, as shown in the table, when comparing the results of the noisy cases under both the LS and MX strategies, our results are better than those of GSSS and SUST, which are also shown in the left bottom and middle of (b-2) - (f-2) and the right bottom of (b-4) - (f-4) in Fig.~\ref{fig:results_PWC_PRE}. 
This indicates the effectiveness of our method in the presence of noise.

We also compare the sampling operators designed by the proposed method.
As an example, Fig.~\ref{fig:results_sampM} visualizes the sampling operators designed by our method for the graph signal generated under the stochastic prior in the experiment.
In the Designs (ii) and (iii), the values of each element are visualized on a log scale in order to highlight the differences in the values.

In the Design (i), the sampling operator is designed to have the Frobenius norm of $\sampM$ less than or equal to $\paramLBall$, which resulted in the elements of $\sampM$ being set between about $-1.5$ and $1.5$ in the visualized case, indicating that the original signal values are used relatively universally to generate the sampled signal.
In the Design (ii), each element of the sampling operator is designed in the range $[0,1]$, and each element of the sampled signal is constructed such that the signal values of some representative vertex are mixed with the weighted signal values of several other vertices.
As in the Design (ii), each element of the sampling operator is designed in the range of $[0,1]$ in the Design (iii); however, the design is sparse, with most elements taking either $0$ or $1$, i.e., the original signal values are uniquely selected or discarded and mixed to construct each element of the sampled signal.

\section{Conclusion}
\label{sec:conclusion}
In this paper, we addressed the challenge of designing a flexible sampling operator for sampling beyond bandlimited graph signals under the arbitrary priors in the vertex domain.
We formulated designing a flexible sampling operator as a problem with an invertibility constraint to achieve the best possible recovery based on the generalized sampling theory.
To handle the constraint, we transformed the problem into a DC optimization problem by relaxing the constraint by using a nuclear norm.
To solve the problem, we developed an efficient solver based on GDPGDC algorithm, which ensures the convergence to a critical point.

Unlike conventional methods with vertex-wise sampling, our approach introduces the innovative concept of mixing values on multiple vertices to generate a sampled signal through designing a flexible sampling operator.
This innovation provides a level of flexibility and efficiency that is lacking in existing methods.
To evaluate the effectiveness of our method, we performed sampling and recovering experiments on various types of graph signals, comparing the MSE between the original and recovered signals. 
The result demonstrated the effectiveness of our approach and highlighted its value in sampling graph signals under arbitrary priors in the vertex domain through designing a flexible sampling operator.

\appendices
\section{The Caluculation of Eq.~\texorpdfstring{\eqref{eq:proj_FNandBox}}{(46)}}
\label{appendix:calc_FNandBox}
Due to~\eqref{eq:prox_f2}, the proximity operator of $(\paramBox \FN{\cdot}^2 + \iota_{\BoxconNM})$ in Eq.~\eqref{eq:proj_FNandBox} can be transformed for $i = 1, \dots, \numv$ and $j = 1, \dots, \nums$ as follows:
\begin{align}
    \label{eq:proof_FNandBox_1}
    &\prox_{\gamma(\paramBox \FN{\cdot}^2 + \iota_{\BoxconNM})} (\mat) \nonumber \\ 
    &= \argmin_{\matt} \paramBox \FN{\matt}^{2} + \indi{\BoxconNM}{\matt} + \frac{1}{2\gamma} \FN{\mat - \matt}^{2} \nonumber \\
    &= \argmin_{\matt} \sum_{i,j} \left[ \paramBox \elee^{2} + \indi{\Boxcon}{\elee} + \frac{1}{2\gamma} \left(\ele - \elee \right)^{2} \right].
\end{align}
Since there are no interaction terms between different $i, j$, the optimization can be performed independently for each element.
From properties of the convex optimization of single variable functions with box constraints, we have
\begin{align}
    \label{eq:proof_FNandBox_3}
    &\argmin_{\elee} \left[ \paramBox \elee^{2} + \indi{\Boxcon}{\elee} + \frac{1}{2\gamma} (\ele - \elee)^{2} \right] \\
    &= 
    \begin{cases}
        \lbound, \ &\mbox{if} \ \prox_{\gamma \paramBox \FN{\cdot}^2} (\ele) < \lbound; \\
        \prox_{\gamma \paramBox \FN{\cdot}^2} (\ele), \ &\mbox{if} \ \lbound \leq \prox_{\gamma \paramBox \FN{\cdot}^2} (\ele) \leq \ubound; \\
        \ubound, \ &\mbox{if} \ \prox_{\gamma \paramBox \FN{\cdot}^2} (\ele) > \ubound
    \end{cases} \nonumber \\
    &= \max \left\{ \lbound, \min \left\{ \prox_{\gamma \paramBox \FN{\cdot}^2} (\ele), \ubound \right\} \right\},
\end{align}
where $\prox_{\gamma \paramBox \FN{\cdot}^2} (\ele)$ is given by
\begin{align}
    \label{eq:proof_FNandBox_2}
    \prox_{\gamma \paramBox \FN{\cdot}^2} (\ele) = \frac{\ele}{1 + 2 \gamma \paramBox}.
\end{align}
Therefore, $\prox_{\gamma( \paramBox \FN{\cdot}^2 + \iota_{\BoxconNM})} (\mat)$ is calculated as follows: for all $i$ from $1$ to $\numv$ and $j$ from $1$ to $\nums$,
\begin{align}
    \label{eq:proof_FNandBox_4}
    &\left[\prox_{\gamma( \paramBox \FN{\cdot}^2 + \iota_{\BoxconNM})} (\mat)\right]_{ij} \nonumber \\
    &= \max \left\{ \lbound, \min \left\{ \prox_{\gamma \paramBox \FN{\cdot}^2} (\ele), \ubound \right\} \right\}.
\end{align}

\section{The Caluculation of Eq.~\texorpdfstring{\eqref{eq:proj_L1andBox}}{(47)}}
\label{appendix:calc_L1andBox}
Due to~\eqref{eq:prox_f2}, the proximity operator of $(\paramBox \Lone{\cdot} + \indi{\BoxconNM}{\cdot})$ in Eq.~\eqref{eq:proj_L1andBox} can be transformed as follows:
\begin{align}
    \label{eq:proof_L1andBox_1}
    &\prox_{\gamma(\paramBox \Lone{\cdot} + \indi{\BoxconNM}{\cdot})} (\mat) \nonumber \\ 
    &= \argmin_{\matt} \paramBox \Lone{\matt} + \indi{\BoxconNM}{\matt} + \frac{1}{2\gamma} \FN{\mat - \matt}^{2} \nonumber \\
    &= \argmin_{\matt} \sum_{i,j} \left[ \paramBox | \elee | + \indi{\Boxcon}{\elee} + \frac{1}{2\gamma} \left(\ele - \elee \right)^{2} \right].
\end{align}
Since there are no interaction terms between different $i$, $j$, the optimization can be performed independently for each element.
From properties of the convex optimization of single variable functions with box constraints, we have
\begin{align}
    \label{eq:proof_L1andBox_2}
    &\argmin_{\elee} \left[ \paramBox | \elee | + \indi{\Boxcon}{\elee} + \frac{1}{2\gamma} \left(\ele - \elee \right)^{2} \right] \nonumber \\
    &= \begin{cases}
        \lbound, 
            &\mathrm{if} \ \prox_{\gamma \paramBox \Lone{\cdot}}\left(\ele\right) < \lbound; \\
        \prox_{\gamma \paramBox \Lone{\cdot}}\left(\ele\right),
            &\mathrm{if} \ \lbound \leq \prox_{\gamma \paramBox \Lone{\cdot}}\left(\ele\right) \leq \ubound; \\
        \ubound,
            &\mathrm{if} \ \prox_{\gamma \paramBox \Lone{\cdot}}\left(\ele\right) > \ubound
    \end{cases} \nonumber \\
    &= \max \left\{ \lbound, \min\left\{ \prox_{\gamma \paramBox \Lone{\cdot}}\left(\ele\right), \ubound \right\}\right\},
\end{align}
where $\prox_{\gamma \paramBox \Lone{\cdot}}\left(\ele\right)$ is given by
\begin{equation}
\label{eq:proof_L1andBox_3}
    \begin{aligned}
        &\prox_{\gamma \paramBox \Lone{\cdot}}\left(\ele\right) \\
        &=
        \begin{cases}
            \ele + \gamma \paramBox, \ &\mbox{if} \ \ele < - \gamma \paramBox; \\
            0, \ &\mbox{if} \ -\gamma \paramBox \leq \ele \leq \gamma \paramBox; \\
            \ele - \gamma \paramBox, \ &\mbox{if} \ \ele > \gamma \paramBox.
        \end{cases}
    \end{aligned}
\end{equation}
Therefore, $\prox_{\gamma(\paramBox \Lone{\cdot} + \iota_{\BoxconNM})} (\mat)$ is calculated as follows: for all $i$ from $1$ to $\numv$ and $j$ from $1$ to $\nums$,
\begin{align}
    &\left[\prox_{\gamma( \paramBox \Lone{\cdot} + \iota_{\BoxconNM})} (\mat)\right]_{ij} \nonumber \\
    &= \max \left\{ \lbound, \min\left\{ \prox_{\gamma \paramBox \Lone{\cdot}}\left(\ele\right), \ubound \right\}\right\}.
\end{align}

\ifCLASSOPTIONcaptionsoff
  \newpage
\fi

\bibliographystyle{IEEEbib}
\begin{thebibliography}{10}

    \bibitem{narang2012perfect}
    S.~K. Narang and A.~Ortega,
    \newblock ``Perfect reconstruction two-channel wavelet filter banks for graph structured data,''
    \newblock {\em IEEE Trans. Signal Process.}, vol. 60, no. 6, pp. 2786--2799, 2012.
    
    \bibitem{agaskar2013spectral}
    A.~Agaskar and Y.~M. Lu,
    \newblock ``A spectral graph uncertainty principle,''
    \newblock {\em IEEE Trans. Inf. Theory}, vol. 59, no. 7, pp. 4338--4356, 2013.
    
    \bibitem{sandryhaila2014big}
    A.~Sandryhaila and J.~M.~F. Moura,
    \newblock ``Big data analysis with signal processing on graphs: Representation and processing of massive data sets with irregular structure,''
    \newblock {\em IEEE Signal Process. Mag.}, vol. 31, no. 5, pp. 80--90, 2014.
    
    \bibitem{shuman2015multiscale}
    D.~I. Shuman, M.~J. Faraji, and P.~Vandergheynst,
    \newblock ``A multiscale pyramid transform for graph signals,''
    \newblock {\em IEEE Trans. Signal Process.}, vol. 64, no. 8, pp. 2119--2134, 2015.
    
    \bibitem{tanaka2020gensamp}
    Y.~Tanaka and Y.~C. Eldar,
    \newblock ``Generalized sampling on graphs with subspace and smoothness priors,''
    \newblock {\em IEEE Trans. Signal Process.}, vol. 68, pp. 2272--2286, 2020.
    
    \bibitem{dong2016learning}
    X.~Dong, D.~Thanou, P.~Frossard, and P.~Vandergheynst,
    \newblock ``Learning {L}aplacian matrix in smooth graph signal representations,''
    \newblock {\em IEEE Trans. Signal Process.}, vol. 64, no. 23, pp. 6160--6173, 2016.
    
    \bibitem{egilmez2018graph}
    H.~E. Egilmez, E.~Pavez, and A.~Ortega,
    \newblock ``Graph learning from filtered signals: Graph system and diffusion kernel identification,''
    \newblock {\em IEEE Trans. Signal Inf. Process. Netw.}, vol. 5, no. 2, pp. 360--374, 2018.
    
    \bibitem{liu2019graph}
    Y.~Liu, L.~Yang, K.~You, W.~Guo, and W.~Wang,
    \newblock ``Graph learning based on spatiotemporal smoothness for time-varying graph signal,''
    \newblock {\em IEEE Access}, vol. 7, pp. 62372--62386, 2019.
    
    \bibitem{dong2019learning}
    X.~Dong, D.~Thanou, M.~Rabbat, and P.~Frossard,
    \newblock ``Learning graphs from data: A signal representation perspective,''
    \newblock {\em IEEE Signal Process. Mag.}, vol. 36, no. 3, pp. 44--63, 2019.
    
    \bibitem{ono2015total}
    S.~Ono, I.~Yamada, and I.~Kumazawa,
    \newblock ``Total generalized variation for graph signals,''
    \newblock in {\em Proc. IEEE Int. Conf. Acoust., Speech, Signal Process. (ICASSP)}, 2015, pp. 5456--5460.
    
    \bibitem{onuki2016graph}
    M.~Onuki, S.~Ono, M.~Yamagishi, and Y.~Tanaka,
    \newblock ``Graph signal denoising via trilateral filter on graph spectral domain,''
    \newblock {\em IEEE Trans. Signal Inf. Process. Netw.}, vol. 2, no. 2, pp. 137--148, 2016.
    
    \bibitem{nagahama2022graph}
    M.~Nagahama, K.~Yamada, Y.~Tanaka, S.~H. Chan, and Y.~C. Eldar,
    \newblock ``Graph signal restoration using nested deep algorithm unrolling,''
    \newblock {\em IEEE Trans. Signal Process.}, vol. 70, pp. 3296--3311, 2022.
    
    \bibitem{berger2017graph}
    Peter Berger, Gabor Hannak, and Gerald Matz,
    \newblock ``Graph signal recovery via primal-dual algorithms for total variation minimization,''
    \newblock {\em IEEE J. Sel. Top. Signal Process}, vol. 11, no. 6, pp. 842--855, 2017.
    
    \bibitem{yamagata2025robust}
    E.~Yamagata, K.~Naganuma, and S.~Ono,
    \newblock ``Robust time-varying graph signal recovery for dynamic physical sensor network data,''
    \newblock {\em IEEE Trans. Signal Inf. Process. Netw.}, 2025.
    
    \bibitem{cheung2018graph}
    G.~Cheung, E.~Magli, Y.~Tanaka, and M.~K. Ng,
    \newblock ``Graph spectral image processing,''
    \newblock {\em Proc. IEEE}, vol. 106, no. 5, pp. 907--930, 2018.
    
    \bibitem{wang2019dynamic}
    Y.~Wang, Y.~Sun, Z.~Liu, S.~E. Sarma, M.~M. Bronstein, and J.~M. Solomon,
    \newblock ``Dynamic graph cnn for learning on point clouds,''
    \newblock {\em ACM Trans. Graph.}, vol. 38, no. 5, pp. 1--12, 2019.
    
    \bibitem{zare2020fusing}
    M.~Zare, M.~S. Helfroush, and K.~Kazemi,
    \newblock ``Fusing hyperspectral and multispectral images using smooth graph signal modelling,''
    \newblock {\em Int. J. Remote Sens.}, vol. 41, no. 22, pp. 8610--8630, 2020.
    
    \bibitem{takemoto2022graph}
    S.~Takemoto, K.~Naganuma, and S.~Ono,
    \newblock ``Graph spatio-spectral total variation model for hyperspectral image denoising,''
    \newblock {\em IEEE Trans. Geosci. Remote Sens.}, vol. 19, pp. 1--5, 2022.
    
    \bibitem{bronstein2017geometric}
    M.~M. Bronstein, J.~Bruna, Y.~LeCun, A.~Szlam, and P.~Vandergheynst,
    \newblock ``Geometric deep learning: going beyond euclidean data,''
    \newblock {\em IEEE Signal Process. Mag.}, vol. 34, no. 4, pp. 18--42, 2017.
    
    \bibitem{chen2014semi}
    S.~Chen, F.~Cerda, P.~Rizzo, J.~Bielak, J.~H. Garrett, and J.~Kova{\v{c}}evi{\'c},
    \newblock ``Semi-supervised multiresolution classification using adaptive graph filtering with application to indirect bridge structural health monitoring,''
    \newblock {\em IEEE Trans. Signal Process.}, vol. 62, no. 11, pp. 2879--2893, 2014.
    
    \bibitem{gadde2014active}
    A.~Gadde, A.~Anis, and A.~Ortega,
    \newblock ``Active semi-supervised learning using sampling theory for graph signals,''
    \newblock in {\em Proc. 20th ACM SIGKDD Int. Conf. Knowl. Discov. Data Mining}, 2014, pp. 492--501.
    
    \bibitem{perozzi2014deepwalk}
    B.~Perozzi, R.~Al-Rfou, and S.~Skiena,
    \newblock ``Deepwalk: Online learning of social representations,''
    \newblock in {\em Proc. 20th ACM SIGKDD Int. Conf. Knowl. Discov. Data Mining}, 2014, pp. 701--710.
    
    \bibitem{saboksayr2021learning}
    S.~S. Saboksayr, G.~Mateos, and M.~Cetin,
    \newblock ``Online graph learning under smoothness priors,''
    \newblock in {\em Proc. Eur. Signal Process. Conf. (EUSIPCO)}, 2021, pp. 1820--1824.
    
    \bibitem{hasanzadeh2019piecewise}
    A.~Hasanzadeh, X.~Liu, N.~Duffield, and K.~R. Narayanan,
    \newblock ``Piecewise stationary modeling of random processes over graphs with an application to traffic prediction,''
    \newblock in {\em Proc. IEEE Int. Conf. Big Data}. IEEE, 2019, pp. 3779--3788.
    
    \bibitem{shuman2013emerging}
    D.~I. Shuman, S.~K. Narang, P.~Frossard, A.~Ortega, and P.~Vandergheynst,
    \newblock ``The emerging field of signal processing on graphs: Extending high-dimensional data analysis to networks and other irregular domains,''
    \newblock {\em IEEE Signal Process. Mag.}, vol. 30, no. 3, pp. 83--98, 2013.
    
    \bibitem{sandryhaila2013discrete}
    A.~Sandryhaila and J.~M.~F. Moura,
    \newblock ``Discrete signal processing on graphs,''
    \newblock {\em IEEE Trans. Signal Process.}, vol. 61, no. 7, pp. 1644--1656, 2013.
    
    \bibitem{ortega2018graph}
    A.~Ortega, P.~Frossard, J.~Kovačević, J.~M.~F. Moura, and P.~Vandergheynst,
    \newblock ``Graph signal processing: Overview, challenges, and applications,''
    \newblock {\em Proceedings of the IEEE}, vol. 106, no. 5, pp. 808--828, 2018.
    
    \bibitem{chen2015discrete}
    S.~Chen, R.~Varma, A.~Sandryhaila, and J.~Kovačević,
    \newblock ``Discrete signal processing on graphs: Sampling theory,''
    \newblock {\em IEEE Trans. Signal Process.}, vol. 63, no. 24, pp. 6510--6523, 2015.
    
    \bibitem{tanaka2020sampling}
    Y.~Tanaka, Y.~C. Eldar, A.~Ortega, and G.~Cheung,
    \newblock ``Sampling signals on graphs: From theory to applications,''
    \newblock {\em IEEE Signal Process. Mag.}, vol. 37, no. 6, pp. 14--30, 2020.
    
    \bibitem{marques2015sampling}
    A.~G. Marques, S.~Segarra, G.~Leus, and A.~Ribeiro,
    \newblock ``Sampling of graph signals with successive local aggregations,''
    \newblock {\em IEEE Trans. Signal Process.}, vol. 64, no. 7, pp. 1832--1843, 2015.
    
    \bibitem{anis2016efficient}
    A.~Anis, A.~Gadde, and A.~Ortega,
    \newblock ``Efficient sampling set selection for bandlimited graph signals using graph spectral proxies,''
    \newblock {\em IEEE Trans. Signal Process.}, vol. 64, no. 14, pp. 3775--3789, 2016.
    
    \bibitem{tsitsvero2016signals}
    M.~Tsitsvero, S.~Barbarossa, and P.~Di~Lorenzo,
    \newblock ``Signals on graphs: Uncertainty principle and sampling,''
    \newblock {\em IEEE Trans. Signal Process.}, vol. 64, no. 18, pp. 4845--4860, 2016.
    
    \bibitem{valsesia2018sampling}
    D.~Valsesia, G.~Fracastoro, and E.~Magli,
    \newblock ``Sampling of graph signals via randomized local aggregations,''
    \newblock {\em IEEE Trans. Signal Inf. Process. Netw.}, vol. 5, no. 2, pp. 348--359, 2018.
    
    \bibitem{tanaka2018spectral}
    Y.~Tanaka,
    \newblock ``Spectral domain sampling of graph signals,''
    \newblock {\em IEEE Trans. Signal Process.}, vol. 66, no. 14, pp. 3752--3767, 2018.
    
    \bibitem{puy2018random}
    G.~Puy, N.~Tremblay, R.~Gribonval, and P.~Vandergheynst,
    \newblock ``Random sampling of bandlimited signals on graphs,''
    \newblock {\em Appl. Comput. Harmon. Anal.}, vol. 44, no. 2, pp. 446--475, 2018.
    
    \bibitem{sakiyama2019eigendecomposition}
    A.~Sakiyama, Y.~Tanaka, T.~Tanaka, and A.~Ortega,
    \newblock ``Eigendecomposition-free sampling set selection for graph signals,''
    \newblock {\em IEEE Trans. Signal Process.}, vol. 67, no. 10, pp. 2679--2692, 2019.
    
    \bibitem{bai2020fast}
    Y.~Bai, F.~Wang, G.~Cheung, Y.~Nakatsukasa, and W.~Gao,
    \newblock ``Fast graph sampling set selection using gershgorin disc alignment,''
    \newblock {\em IEEE Trans. Signal Process.}, vol. 68, pp. 2419--2434, 2020.
    
    \bibitem{jayawant2022practical}
    A.~Jayawant and A.~Ortega,
    \newblock ``Practical graph signal sampling with log-linear size scaling,''
    \newblock {\em Signal Process.}, vol. 194, pp. 108436, 2022.
    
    \bibitem{li2025sensor}
    J.~Li and B.~Cai,
    \newblock ``Sensor placement method for water distribution networks based on sampling of non-bandlimited graph signals,''
    \newblock {\em Digital Signal Processing}, vol. 156, pp. 104809, 2025.
    
    \bibitem{elder2009beyond}
    Y.~C. Eldar and T.~Michaeli,
    \newblock ``Beyond bandlimited sampling,''
    \newblock {\em IEEE Signal Process. Mag.}, vol. 26, no. 3, pp. 48--68, 2009.
    
    \bibitem{eldar2015sampling}
    Y.~C. Eldar,
    \newblock {\em Sampling theory: Beyond bandlimited systems},
    \newblock Cambridge University Press, 2015.
    
    \bibitem{chepuri2018graphsamp}
    S.~P. Chepuri, Y.~C. Eldar, and G.~Leus,
    \newblock ``Graph sampling with and without input priors,''
    \newblock in {\em Proc. IEEE Int. Conf. Acoust., Speech, Signal Process. (ICASSP)}, 2018, pp. 4564--4568.
    
    \bibitem{hara2020generalized}
    J.~Hara, Y.~Tanaka, and Y.~C. Eldar,
    \newblock ``Generalized graph spectral sampling with stochastic priors,''
    \newblock in {\em Proc. IEEE Int. Conf. Acoust., Speech, Signal Process. (ICASSP)}. IEEE, 2020, pp. 5680--5684.
    
    \bibitem{hara2021design}
    J.~Hara, K.~Yamada, S.~Ono, and Y.~Tanaka,
    \newblock ``Design of graph signal sampling matrices for arbitrary signal subspaces,''
    \newblock in {\em Proc. IEEE Int. Conf. Acoust., Speech, Signal Process. (ICASSP)}, 2021, pp. 5275--5279.
    
    \bibitem{hara2022gsss}
    J.~Hara and Y.~Tanaka,
    \newblock ``Sampling set selection for graph signals under arbitrary signal priors,''
    \newblock in {\em Proc. IEEE Int. Conf. Acoust., Speech, Signal Process. (ICASSP)}, 2022, pp. 5732--5736.
    
    \bibitem{hara2023graph}
    J.~Hara, Y.~Tanaka, and Y.~C. Eldar,
    \newblock ``Graph signal sampling under stochastic priors,''
    \newblock {\em IEEE Trans. Signal Process.}, vol. 71, pp. 1421--1434, 2023.
    
    \bibitem{banert2019general}
    S.~Banert and R.~I. Boț,
    \newblock ``A general double-proximal gradient algorithm for dc programming,''
    \newblock {\em Math. Program.}, vol. 178, no. 1-2, pp. 301--326, 2019.
    
    \bibitem{yamashita2024apsipa}
    K.~Yamashita, K.~Naganuma, and S.~Ono,
    \newblock ``Generalized graph signal sampling under subspace priors by difference-of-convex minimization,''
    \newblock in {\em Proc. Asia-Pac. Signal Inf. Process. Assoc. Annu. Summit Conf. (APSIPA ASC)}, 2024, pp. 1--6.
    
    \bibitem{eldar2003sampling}
    Y.~C. Eldar,
    \newblock ``Sampling with arbitrary sampling and reconstruction spaces and oblique dual frame vectors,''
    \newblock {\em Journal of Fourier Analysis and Applications}, vol. 9, pp. 77--96, 2003.
    
    \bibitem{bauschke2011convex}
    H.~H. Bauschke and P.~L. Combettes,
    \newblock {\em Convex analysis and monotone operator theory in Hilbert spaces},
    \newblock Springer, 2011.
    
    \bibitem{bauschke2017correction}
    H.~H. Bauschke and P.~L. Combettes,
    \newblock {\em Correction to: convex analysis and monotone operator theory in Hilbert spaces},
    \newblock Springer, 2017.
    
    \bibitem{fazel2002matrix}
    M.~Fazel,
    \newblock {\em Matrix rank minimization with applications},
    \newblock Ph.D. thesis, PhD thesis, Stanford University, 2002.
    
    \bibitem{yamashita2025icassp}
    K.~Yamashita, K.~Naganuma, and S.~Ono,
    \newblock ``Controlling the number of sample-contributive vertices in generalized sampling of graph signals,''
    \newblock in {\em Proc. IEEE Int. Conf. Acoust., Speech, Signal Process. (ICASSP)}, 2025,
    \newblock (Accepted).
    
    \bibitem{perraudin2014gspbox}
    N.~Perraudin, J.~Paratte, D.~Shuman, L.~Martin, V.~Kalofolias, P.~Vandergheynst, and D.~K. Hammond,
    \newblock ``{GSPBOX}: A toolbox for signal processing on graphs,''
    \newblock {\em arXiv}, 2014,
    \newblock [Online]. Available: https://arxiv.org/abs/1408.5781.
    
    \bibitem{chen2016representations}
    S.~Chen, R.~Varma, A.~Singh, and J.~Kova{\v{c}}evi{\'c},
    \newblock ``Representations of piecewise smooth signals on graphs,''
    \newblock in {\em Proc. IEEE Int. Conf. Acoust., Speech, Signal Process. (ICASSP)}. IEEE, 2016, pp. 6370--6374.
    
    \bibitem{gadde2015probabilistic}
    A.~Gadde and A.~Ortega,
    \newblock ``A probabilistic interpretation of sampling theory of graph signals,''
    \newblock in {\em Proc. IEEE Int. Conf. Acoust., Speech, Signal Process. (ICASSP)}, 2015, pp. 3257--3261.
    
    \end{thebibliography}
    

\begin{IEEEbiography}[{\includegraphics[width=1in,height=1.25in,clip,keepaspectratio]{author/yamashita.jpg}}]{Keitaro Yamashita}
    (Graduate Student Member, IEEE) received the B.E. degree in control engineering from Tokyo Institute of Technology, Japan, in 2015, and the M.E. degree in decision science and technology from Tokyo Institute of Technology, Japan, in 2018. 
    After receiving the M.E. degree, he joined the Ministry of Economy, Trade and Industry (METI), Japan, and was in charge of the policy planning for the research and development, deep tech startups, the carbon neutrality, and the data free flow with trust (DFFT).
    He has also worked for the Agency for Cultural Affairs (ACA), Japan, where he was in charge of the copyright policy.
    He is currently pursuing the Ph.D. degree in the Department of Computer Science at Institute of Science Tokyo.
    He recieved the Best Paper Award at Asia Pacific Signal and Information Processing Association Annual Summit and Conference (APSIPA ASC) in 2024.
    His research interests include graph signal processing and optimization.
\end{IEEEbiography}

\begin{IEEEbiography}[{\includegraphics[width=1in,height=1.25in,clip,keepaspectratio]{author/naganuma.jpg}}]{Kazuki Naganuma}
    (Member, IEEE) received the B.E. degree in information and computer sciences from Kanagawa Institute of Technology, Atsugi, Japan, in 2020, and the M.E. degree and Ph.D. degree in information and computer sciences from Tokyo Institute of Technology, Yokohama, Japan, in 2022, where he is currently a researcher at Institute of Science Tokyo. Since October 2023, he has been a Research Fellow (DC2) of Japan Society for the Promotion of Science (JSPS) and a Researcher of ACT-X of Japan Science and Technology Agency (JST), Tokyo, Japan. His research interests include signal and image processing and optimization theory. Dr. Naganuma received the Student Conference Paper Award from the IEEE SPS Japan Chapter in 2023.
\end{IEEEbiography}

\begin{IEEEbiography}[{\includegraphics[width=1in,height=1.25in,clip,keepaspectratio]{author/ono.jpg}}]{Shunsuke Ono}
    Shunsuke Ono (S’11–M’15–SM'23) received a B.E. degree in Computer Science in 2010 and M.E. and Ph.D. degrees in Communications and Computer Engineering in 2012 and 2014 from the Tokyo Institute of Technology, respectively.

    From April 2012 to September 2014, he was a Research Fellow (DC1) of the Japan Society for the Promotion of Science (JSPS). He is currently an Associate Professor in the Department of Computer Science, School of Computing, Tokyo Institute of Technology. From October 2016 to March 2020 and from October 2021 to present, he was/is a Researcher of Precursory Research for Embryonic Science and Technology (PRESTO), Japan Science and Technology Agency (JST), Tokyo, Japan. His research interests include signal processing, image analysis, remote sensing, mathematical optimization, and data science.
    
    Dr. Ono received the Young Researchers’ Award and the Excellent Paper Award from the IEICE in 2013 and 2014, respectively, the Outstanding Student Journal Paper Award and the Young Author Best Paper Award from the IEEE SPS Japan Chapter in 2014 and 2020, respectively, the Funai Research Award from the Funai Foundation in 2017, the Ando Incentive Prize from the Foundation of Ando Laboratory in 2021, the Young Scientists’ Award from MEXT in 2022, and the Outstanding Editorial Board Member Award from IEEE SPS in 2023. He has been an Associate Editor of IEEE TRANSACTIONS ON SIGNAL AND INFORMATION PROCESSING OVER NETWORKS since 2019.
\end{IEEEbiography}

\end{document}